\documentclass[12pt,a4paper]{article}

\input{glyphtounicode}
\pdfgentounicode=1

\usepackage[utf8]{inputenc}
\usepackage[T1]{fontenc}
\usepackage{lmodern}
\usepackage[ngerman]{babel}
\usepackage{amsmath,amssymb,amsthm,mathtools}
\usepackage{xurl}
\usepackage[colorlinks=true,linkcolor=blue,citecolor=red,urlcolor=blue,unicode=true]{hyperref}
\usepackage{geometry}
\geometry{margin=2.5cm}
\usepackage{titling}
\usepackage{enumitem}
\usepackage{array}
\usepackage{booktabs}
\usepackage{tabularx}
\usepackage[capitalise,noabbrev]{cleveref}

\newtheorem{theorem}{Theorem}[section]
\newtheorem{lemma}[theorem]{Lemma}
\newtheorem{proposition}[theorem]{Proposition}
\newtheorem{corollary}[theorem]{Korollar}
\newtheorem{conjecture}[theorem]{Vermutung}
\theoremstyle{definition}
\newtheorem{definition}[theorem]{Definition}
\theoremstyle{remark}
\newtheorem{remark}[theorem]{Bemerkung}

\newcommand{\HH}{\mathbb{H}}
\newcommand{\RR}{\mathbb{R}}
\newcommand{\CC}{\mathbb{C}}
\newcommand{\NN}{\mathbb{N}}
\newcommand{\ZZ}{\mathbb{Z}}
\newcommand{\QW}{\mathrm{QW}}
\newcommand{\Sel}{Z_{\Gamma}}
\newcommand{\Delop}{\Delta_{\Gamma}}
\newcolumntype{L}[1]{>{\raggedright\arraybackslash}p{#1}}
\newcolumntype{Y}{>{\raggedright\arraybackslash}X}
\setlength{\droptitle}{-3.1em}
\pretitle{\begin{center}\Large}
\posttitle{\par\end{center}\vspace{-0.8em}}
\preauthor{\begin{center}\normalsize}
\postauthor{\par\end{center}\vspace{-0.7em}}
\predate{\begin{center}\normalsize}
\postdate{\par\end{center}\vspace{-1.2em}}

\title{Das Scheitern von NE-B als Hilbert--Pólya-Detektion:\\
Universalität der Weilschen quadratischen Form v2.0 auf die Selbergsche Zetafunktion}
\author{Lukas Geiger\\
\small Bernau, Deutschland\\
\small ORCID: \href{https://orcid.org/0009-0005-7296-1534}{0009-0005-7296-1534}}
\date{\today}

\hypersetup{
  pdftitle={Das Scheitern von NE-B als Hilbert--Pólya-Detektion: Universalität der Weilschen quadratischen Form v2.0 auf die Selbergsche Zetafunktion},
  pdfauthor={Lukas Geiger},
  pdfsubject={Selberg v2.0 Transferpapier zur Weilschen quadratischen Form},
  pdfkeywords={Selbergsche Zetafunktion, Weilsche quadratische Form, Verschiebungs-Paritäts-Lemma, Hilbert--Pólya, hyperbolische Fläche, Laplace-Operator, NE-B}
}
\makeatletter
\renewcommand*{\HyperDestNameFilter}[1]{de.#1}
\makeatother

\begin{document}
\overfullrule=0pt
\setlength{\emergencystretch}{3em}
\maketitle
\vspace{-2.0em}

\begin{abstract}
Die k\"urzlich f\"ur die Riemannsche Vermutung (Geiger 2026,
Teil~II) vorgeschlagene v2.0-Methodik reduziert die RH auf die \emph{gerade Dominanz} (even dominance) der Weilschen quadratischen Form
$\QW_\lambda$ mittels dreier algebraischer und analytischer Bestandteile: das Verschiebungs-Parit\"ats-Lemma
(Shift Parity Lemma), die Dominanz der Frontier-Primzahlen und zwei Nichtexistenz-Theoreme (NE-A:
Nicht-Positivit\"at des Primzahl-Verschiebungs-Multiplikators $M(\xi) = -2\,\mathrm{Re}[\zeta'/\zeta(\tfrac{1}{2}+i\xi)]$,
NE-B: kein universeller kommutierender Operator). Eine nat\"urliche Frage ist, ob dieser
Rahmen spezifisch f\"ur die Riemannsche Zetafunktion ist oder ob es sich um eine
allgemeine Methode f\"ur Zeta-Strukturen mit expliziten Formeln handelt. Wir testen sie an der
Selbergschen Zetafunktion $\Sel(s)$ einer kompakten hyperbolischen Fl\"ache
$\Gamma\backslash\HH$. Wir zeigen:

\begin{enumerate}[label=(\roman*),leftmargin=*,itemsep=0.15em,topsep=0.25em,parsep=0pt,partopsep=0pt]
\item Das Verschiebungs-Parit\"ats-Lemma \"ubertr\"agt sich unver\"andert (es ist rein algebraisch).
\item Die Frontier-Dominanz \"ubertr\"agt sich mit modifizierten Konstanten, was der
exponentiellen Dichte primitiver geschlossener Geod\"aten $\pi_\Gamma(T) \sim e^T/T$
anstelle von $\pi(x) \sim x/\log x$ Rechnung tr\"agt.
\item Ein Analogon $\QW^\Gamma_\lambda$ der Weilschen quadratischen Form ist auf
der Testfunktionsseite der Selbergschen Spurformel wohldefiniert; die Beziehung
zum geometrischen Fluss wird durch die Selberg-Transformation vermittelt, nicht
durch einen hier konstruierten unit\"aren Intertwiner.
\item \emph{Die gerade Dominanz von $\QW^\Gamma_\lambda$} liefert zusammen mit
der Bedingung ohne Ausnahmeeigenwerte $\lambda_1(\Delop)\ge 1/4$ und einem
expliziten $C2$/Minimizer-Identifikationsgate eine bedingte
v2.0-Konsistenzreduktion der strengen Aussage für die nicht-trivialen
spektralen Nullstellen im kritischen Streifen; dies
ist ein Konsistenztest gegen die klassische spektrale Beschreibung, kein neuer
unbedingter Beweis.
\item \emph{NE-B ist im sph\"arischen Selberg-Kanal nicht vorhanden.} Der
Casimir-/Laplace-Operator $\Delop$ kommutiert mit der Algebra radialer
sph\"arischer Faltungsoperatoren und liefert dort den klassischen
Hilbert--P\'olya-Kanal. Individuelle trunkierte Geod\"atenfluss-Pullbacks auf
$SM$ werden damit nicht als Intervall-Verschiebungen identifiziert.
\end{enumerate}

Die Abwesenheit des einschl\"agigen NE-B-Hindernisses im Selberg-Setting ist kein
Mangel; sie spiegelt die geometrische Koordination der Spurformel auf einem
gemeinsamen symmetrischen Raum wider, die bei den Primzahlen fehlt. Wir leiten
ein vorsichtiges Meta-Prinzip ab: der v2.0-Rahmen liefert eine gemeinsame
Quadratform-Testsprache, w\"ahrend ein Hilbert--P\'olya-Operator zus\"atzliche
operatorielle Struktur verlangt. Im Selberg-Fall ist diese Struktur im
sph\"arischen Casimir-Kanal vorhanden; im Riemann-Fall wird sie durch den
getesteten NE-B-Weg nicht geliefert. Die Weilsche quadratische Form, nicht der
Spektraloperator, ist das gemeinsame Objekt.
\end{abstract}
\vspace{-0.8em}

\noindent\textbf{MSC 2020:} 11M36, 11F72, 58J50, 47A75\\
\textbf{Schlüsselwörter:} Selbergsche Zetafunktion, Weilsche quadratische Form, Verschiebungs-Paritäts-Lemma, Hilbert--Pólya, hyperbolische Fläche, Laplace-Operator, NE-B.

\clearpage
\tableofcontents
\clearpage

% =============================================================================
\section{Einleitung}
\label{sec:intro}
% =============================================================================

\subsection{Motivation und Hauptaussage}

Das v2.0-Programm f\"ur die Riemannsche Vermutung \cite{Geiger2026partII}
reduziert die RH auf eine Ungleichung quadratischer Formen. Konkret dr\"uckt Connes' 2026-Artikel
\cite{Connes2026} die RH als die Aussage aus, dass der minimale Eigenwert einer
einparametrigen Familie von quadratischen Formen $\QW_\lambda$ auf
$L^2([-\log\lambda,\log\lambda])$ einfach ist und eine gerade Eigenfunktion besitzt.
Der v2.0-Ansatz beweist dies in drei Schichten:

\begin{itemize}
\item ein rein algebraisches \emph{Verschiebungs-Parit\"ats-Lemma} (Theorem~4.1 in \cite{Geiger2026partII}),
das zeigt, dass die Gerade-Ungerade-Differenzmatrix $D_N(r)$ f\"ur alle $r \in (0,2)$ und $N \geq 3$ einen strikt negativen
minimalen Eigenwert hat;

\item ein analytischer \emph{Frontier-Dominanz}-Mechanismus: die L\"ucke
$\Delta(\lambda) = \lambda_1^+ - \lambda_1^-$ wird von Primzahlen mit
$\log p / \log\lambda$ nahe an $1$ (``Frontier-Primzahlen'') angetrieben, deren Dichte
durch den Primzahlsatz kontrolliert wird;

\item zwei strukturelle Nichtexistenz-Ergebnisse, die alternative Routen ausschlie\ss{}en:
\emph{NE-A} (der Primzahl-Verschiebungs-Multiplikator $M(\xi)$ ist nicht fast \"uberall nichtnegativ,
was totale Positivit\"at $\mathrm{PF}_\infty$ ausschlie\ss{}t) und \emph{NE-B} (es gibt keine nichttriviale
symmetrische Matrix, die mit allen $D_N(r)$ gleichzeitig kommutiert, was
universelle Operatoren vom Sturm--Liouville-Typ bei den getesteten Trunkierungen ausschlie\ss{}t).
\end{itemize}

Das v2.0-Programm ist erfolgreich, \emph{ohne} einen Hilbert--P\'olya-Operator
zu konstruieren. Es ist naheliegend zu fragen, ob dieses methodologische Paket auf
andere zeta-\"ahnliche Funktionen \"ubertragbar ist und, wenn ja, was es uns lehrt. Der paradigmatische
Fall ist die Selbergsche Zetafunktion $\Sel(s)$ einer kompakten hyperbolischen Fl\"ache
$M = \Gamma\backslash\HH$, wobei:

\begin{itemize}
\item die geometrische explizite Formel (Selbergsche Spurformel \cite{Selberg1956})
die Rolle der Weilschen expliziten Formel \cite{Weil1952} \"ubernimmt;
\item die L\"angen primitiver geschlossener Geod\"aten die Rolle der Primzahlen spielen;
\item die spektrale Nullstellenbeschreibung durch die Selbstadjungiertheit des
Laplace-Operators $\Delop$ ein Theorem ist; die strenge kritische-Linien-Form
erfordert zus\"atzlich die Bedingung ohne Ausnahmeeigenwerte
$\lambda_1(\Delop)\ge 1/4$.
\end{itemize}

Wir stellen also zwei konkrete Fragen:

\medskip
\noindent\textbf{F1.}\\
\emph{\"Ubertragen sich die v2.0-Bestandteile
(Verschiebungs-Parit\"ats-Lemma, Frontier-Dominanz,
Weilsche quadratische Form, NE-A, NE-B)
alle auf das Selberg-Setting?}

\medskip
\noindent\textbf{F2.}\\
\emph{Wenn eines von ihnen sich nicht \"ubertragen l\"asst,
was ist der strukturelle Grund daf\"ur?}

\medskip
Unsere Antworten in K\"urze:

\begin{itemize}
\item Verschiebungs-Parit\"ats-Lemma: \textbf{\"ubertr\"agt sich unver\"andert} (algebraisch).
\item Frontier-Dominanz: \textbf{\"ubertr\"agt sich}
mit modifizierten Konstanten,
was $\pi_\Gamma(T) \sim e^T/T$ widerspiegelt.
\item Weilsche quadratische Form: \textbf{nat\"urliches Analogon $\QW^\Gamma_\lambda$} auf der
Testfunktionsseite der Spurformel; die geometrische Seite wird \"uber den
Selberg-Transform gekoppelt.
\item NE-A: \textbf{\"ubertr\"agt sich} -- nach Lokalisierung der Vorzeichenwechsel in den
trivialen, topologischen und Identit\"atsbeitr\"agen der expliziten Formel und
nicht in den Nullstellen auf der kritischen Linie.
\item NE-B: \textbf{das einschl\"agige Hindernis fehlt im sph\"arischen Kanal}. Ein
universeller Kommutant existiert dort: der Casimir-/Laplace-Operator $\Delop$
kommutiert mit radialen sph\"arischen Faltungsoperatoren. Dies ist nicht als
Aussage \"uber alle individuellen Geod\"atenfluss-Pullbacks auf $SM$ zu lesen.
\end{itemize}

Die Abwesenheit des relevanten NE-B-Hindernisses ist das wissenschaftlich
interessante Ergebnis. Es ist kein Scheitern von v2.0; es ist eine
\emph{Diagnostik}. Wo ein solcher Kommutant im passenden Funktionskanal existiert,
ist ein operatorieller Hilbert--P\'olya-Weg verf\"ugbar. Wo der getestete
NE-B-Weg keinen nichttrivialen Kommutanten zul\"asst (wie im Riemann-Modell der
Part-II-Trunkierungen), ist zumindest diese universelle Kommutantenroute
blockiert. In beiden F\"allen ist der v2.0-Rahmen als Quadratform-Test
anwendbar.

\begin{remark}[Position innerhalb des Meta-Rahmens]
Dieses Papier fungiert als der \emph{Lie-Ursprungs-Begleiter} (Lie-origin companion) zum arithmetischen
Riemann-Fall \cite{Geiger2026partII} und reiht sich in die Drei-Komponenten-Zerlegung
$\mathrm{RH}_X = \mathrm{GBC}_X + C2_X + \mathrm{Hurwitz}_X$
ein, die im Begleit-Meta-Papier \cite{GeigerMeta2026} entwickelt wurde. Selberg
ist der paradigmatische Fall, in dem der sph\"arische Lie-/Casimir-Kanal die
operatorielle Seite liefert; die vollst\"andige $C2$-Identifikation mit dem
Minimizer bleibt als explizites Gate sichtbar. Riemann ist der Fall, in dem
der Parit\"ats-Teilschritt durch v2.1 geschlossen wird, aber der Eigenvektor-Realisierungs-Teilschritt
offen bleibt. Siehe \Cref{rem:meta-paper-refinement} unten f\"ur
Details.
\end{remark}

\subsection{Haupttheorem (informell)}

\begin{theorem}[v2.0 Selberg-Transfer, informell]
\label{thm:main-informal}
Sei $M = \Gamma\backslash\HH$ eine kompakte hyperbolische Fl\"ache und
$\Sel(s)$ ihre Selbergsche Zetafunktion. Die v2.0-Methodik ist auf die
Selbergsche quadratische Form $\QW^\Gamma_\lambda$ anwendbar und reduziert die
strenge Aussage für die nicht-trivialen spektralen Nullstellen im kritischen
Streifen bedingt auf gerade Dominanz,
$\lambda_1(\Delop)\ge 1/4$ und ein explizites
$C2$/Minimizer-Identifikationsgate. Das Analogon von NE-A gilt. Das relevante
NE-B-Hindernis ist im sph\"arischen Selberg-Kanal abwesend: der
Laplace--Beltrami-/Casimir-Operator $\Delop$ ist ein universeller Kommutant f\"ur
radiale sph\"arische Faltungsoperatoren. Der resultierende Weg ist daher ein
bedingter Konsistenztest und offenbart ein vorsichtiges Meta-Prinzip: die
Weilsche quadratische Form ist die gemeinsame Testsprache, w\"ahrend ein
Hilbert--P\'olya-Operator zus\"atzliche operatorielle Struktur verlangt.
\end{theorem}

\subsection{Struktur der Arbeit}

Abschnitt~\ref{sec:v2-recollect} ruft die v2.0-Bestandteile in Erinnerung.
Abschnitt~\ref{sec:selberg-setting} f\"uhrt die Selbergsche Spurformel und das
geod\"atische Analogon der expliziten Formel ein.
Abschnitt~\ref{sec:v2-to-selberg} f\"uhrt den Transfer des Verschiebungs-Parit\"ats-Lemmas und der
Frontier-Dominanz durch und konstruiert $\QW^\Gamma_\lambda$; eine bedingte
Reduktion der kritischen-Linien-Aussage wird abgeleitet.
Abschnitt~\ref{sec:ne-b-fails} beweist die Abwesenheit des NE-B-Hindernisses im
sph\"arischen Selberg-Kanal und identifiziert den Casimir-Kommutanten $\Delop$.
Abschnitt~\ref{sec:meta} leitet das Meta-Prinzip ab: Hilbert--P\'olya ist ein
operatorieller Spezialfall eines allgemeineren Quadratform-Tests.
Abschnitt~\ref{sec:applications} diskutiert Anwendungen und Ausblick
(verallgemeinerte L-Funktionen, Ruelle-Zeta, Quantenchaos).
Abschnitt~\ref{sec:conclusion} schlie\ss{}t.


% =============================================================================
\section{R\"uckblick auf die v2.0 Methode}
\label{sec:v2-recollect}
% =============================================================================

Wir fassen die vier Schl\"usselbestandteile des v2.0-Ansatzes zur RH kurz zusammen. Eine
ausf\"uhrliche Behandlung findet sich in \cite{Geiger2026partII}.

\subsection{Die Weilsche quadratische Form}

Setze $L = \log\lambda$. Auf $L^2([-L,L])$ definieren wir
\begin{equation}
\label{eq:QW-riemann}
\begin{split}
\langle A_\lambda f\mid f\rangle
&= (\log 4\pi + \gamma)\|f\|^2 \\
&\quad + \int_0^\infty \Big[\langle T_u f + T_{-u} f, f\rangle - 2e^{-u/2}\|f\|^2\Big]
\frac{e^{u/2}}{2\sinh u}\,du \\
&\quad + \sum_p\sum_{m\ge 1}\frac{\log p}{p^{m/2}}
\langle T_{m\log p} f + T_{-m\log p} f, f\rangle,
\end{split}
\end{equation}
wobei $T_s$ der Verschiebungs-Operator $(T_s f)(t) = f(t-s)\mathbf{1}_{[-L,L]}(t-s)$
ist und $\gamma$ die Euler--Mascheroni-Konstante ist. Die quadratische Form $\QW_\lambda$
ist die quadratische Form von $A_\lambda$.  Das wesentliche algebraische Werkzeug ist
Theorem~6.1 von \cite{Connes2026} (gemeinsam mit van Suijlekom bewiesen,
aufbauend auf \cite{ConnesMoscovici2022}):

\begin{theorem}[{\cite[Theorem~6.1]{Connes2026}}]
\label{thm:connes}
Sei $L > 0$, sei $\mathcal{D}$ eine reelle Distribution auf dem Intervall
$[0,L]$, und sei $\tilde{\mathcal{D}}$ die zugeh\"orige gerade Distribution
auf $[-L,L]$.  Angenommen, die quadratische Form mit Schwartz-Kern
$\tilde{\mathcal{D}}(x-y)$ definiert einen nach unten beschr\"ankten selbstadjungierten Operator auf
$L^2\!\bigl([-\tfrac{L}{2},\tfrac{L}{2}]\bigr)$, und das Minimum seines
Spektrums ist ein einfacher, isolierter Eigenwert mit gerader Eigenfunktion~$\eta$.
Dann liegen alle Nullstellen der ganzen Funktion $\tilde\eta(z)$, $z\in\mathbb{C}$
(Fouriertransformierte von\/~$\eta$), auf der reellen Achse.
\end{theorem}

Angewandt auf die Weilsche quadratische Form $\QW_\lambda$ (mit $\mathcal{D}$ als Weil-Distribution
aus der expliziten Formel der Primzahlen), gilt: wenn der minimale Eigenwert
f\"ur alle $\lambda\ge\lambda_0$ einfach ist und eine gerade Eigenfunktion besitzt, und die
zugeh\"origen Eigenvektoren gegen die Riemannsche $\Xi$-Funktion konvergieren, dann stimmen die
Nullstellen von $\tilde\eta$ mit denen von $\Xi$ \"uberein; folglich gilt die RH.  Der
Konvergenzschritt bleibt offen \cite[\S6.6]{Connes2026}.

\begin{remark}[Intervallkonvention]
\label{rem:interval-convention}
Theorem~\ref{thm:connes} platziert den Operator der quadratischen Form auf
$L^2\!\bigl([-\tfrac{L}{2},\tfrac{L}{2}]\bigr)$ in Connes' Notation.
In der vorliegenden Arbeit (ab Abschnitt~\ref{sec:v2-to-selberg})
leben die Testfunktionen auf $L^2([-L,L])$.  Die Anwendung von
Theorem~\ref{thm:connes} erfordert die Identifikation $L_{\mathrm{Thm}} = 2L$,
d.h.\ die halbe L\"ange in Connes' Aussage entspricht der vollen Trunkierungsl\"ange~$L$.
\end{remark}

\begin{remark}[Unabh\"angigkeit von Connes' Konvergenzschritt]
\label{rem:connes-hedge}
Theorem~\ref{thm:connes} stammt aus Connes' 2026 Preprint \cite{Connes2026};
die Konvergenz der Konstruktion gegen die Weilsche Distribution wird als noch offen angemerkt
\cite[\S6.6]{Connes2026}.  In der Selberg-Anwendung
(Abschnitt~\ref{sec:v2-to-selberg}) ist dieser offene Schritt \emph{unbedeutend}: die
spektrale Parametrisierung der Nullstellen ist unabh\"angig durch Selbergs
klassische Theorie \"uber den Laplace-Operator $\Delop$ gesichert
\cite{Selberg1956}; die strenge kritische-Linien-Form erfordert zus\"atzlich
die Bedingung ohne Ausnahmeeigenwerte $\lambda_1(\Delop)\ge 1/4$.
Theorem~\ref{thm:connes} spielt eine strukturelle Konsistenzrolle, keine logisch
notwendige.  Das breitere Programm der Codierung von $L$-Funktionen \"uber spektrale
quadratische Formen wird zus\"atzlich durch Connes' 1999 Ansatz der nichtkommutativen
Geometrie \cite{Connes1999}, Deningers Programm regul\"arer Determinanten
\cite{Deninger1994} und das Bost--Connes-Thermodynamiksystem
\cite{BostConnes1995} gest\"utzt.
\end{remark}

\subsection{Das Verschiebungs-Parit\"ats-Lemma}

Seien $\varphi_n(t) = \cos(n\pi t/L)/\sqrt{L}$ (gerade Basis) und
$\psi_n(t) = \sin((n+1)\pi t/L)/\sqrt{L}$ (ungerade Basis). Wir definieren die
\emph{Differenzmatrix} $D_N(r)$ durch
\begin{equation}
\label{eq:DN}
D_{nm}(r) = \langle\varphi_n, (S_{rL}+S_{-rL})\varphi_m\rangle - \langle\psi_n, (S_{rL}+S_{-rL})\psi_m\rangle.
\end{equation}

\begin{theorem}[Verschiebungs-Parit\"ats-Lemma, \cite{Geiger2026partII}]
\label{thm:SPL}
F\"ur jedes $N \ge 3$ und jedes $r \in (0, 2)$ gilt $\lambda_{\min}(D_N(r)) < 0$.
\end{theorem}

Dies ist eine rein algebraische, geschlossene Aussage \"uber trigonometrische
Matrizen; die Eintr\"age von $D_N(r)$ sind explizite Summen von $\sin$ und $\cos$
mit linearen und konstanten Koeffizienten in $r$.

\subsection{Frontier-Dominanz}

Die kumulative Gerade-Ungerade-L\"ucke $\Delta(\lambda) = \lambda_1^+(\lambda) - \lambda_1^-(\lambda)$
wird durch Primzahlen im \emph{Frontier-Band} $[\lambda/e, \lambda]$ angetrieben: ihre
Anzahl ist $\pi(\lambda) - \pi(\lambda/e) \sim \lambda/\log\lambda$ laut PNT,
und sie sind diejenigen, f\"ur die $\log p / L \in [1-1/\log\lambda, 1]$ gilt. Nach
Theorem~\ref{thm:SPL} liefert jede solche Primzahl einen negativen Summanden der
Ordnung $-c\,\log p / \sqrt{p}$ zur Differenz, wobei die Summe jede
archimedische Korrektur oder Korrektur kleiner Primzahlen dominiert. Numerisch wird in \cite{Geiger2026partII}
gezeigt, dass Frontier-Primzahlen $> 99\%$ der L\"ucke f\"ur alle getesteten
$\lambda \le 1.3\cdot 10^6$ ausmachen.

\subsection{NE-A und NE-B}

Wir definieren den \emph{Primzahl-Verschiebungs-Multiplikator}
\begin{equation}
\label{eq:Mxi}
M(\xi) = -2\,\mathrm{Re}\Big[\frac{\zeta'}{\zeta}\Big(\tfrac{1}{2}+i\xi\Big)\Big].
\end{equation}

\begin{theorem}[NE-A, \cite{Geiger2026partII}]
\label{thm:NEA-riemann}
Die Menge $\{\xi \in \RR : M(\xi) < 0\}$ hat ein positives Lebesgue-Ma\ss{}. Insbesondere
ist der Primzahl-Verschiebungs-Operator nicht total positiv im Sinne von $\mathrm{PF}_\infty$.
\end{theorem}

\begin{theorem}[NE-B, \cite{Geiger2026partII}]
\label{thm:NEB-riemann}
F\"ur $N \in \{3,5,7,10,15\}$ und eine beliebig hinreichend dichte Stichprobe von $r \in (0,2)$
ist die einzige symmetrische Matrix $T \in \mathrm{Sym}_N(\RR)$, die
$[T, D_N(r_i)] = 0$ f\"ur alle $i$ erf\"ullt, $T = cI$. Folglich existiert bei diesen Trunkierungen
kein nichttrivialer universeller kommutierender Operator; die Gesamtraum-Vermutung
\cite[Conjecture~5.3]{Geiger2026partII} behauptet dasselbe f\"ur
alle $N$.
\end{theorem}

NE-B ist das technisch entscheidende strukturelle Ergebnis: es dr\"uckt aus, dass die
Primzahlen keine verborgene gemeinsame Symmetrie besitzen. Der konzeptionelle Inhalt ist
die multiplikative Unabh\"angigkeit.


% =============================================================================
\section{Das Selberg-Setting}
\label{sec:selberg-setting}
% =============================================================================

Wir skizzieren kurz das Standard-Setup.

\subsection{Kompakte hyperbolische Fl\"achen}

Sei $M = \Gamma\backslash\HH$, wobei $\HH = \{x+iy : y > 0\}$ die obere
Halbebene mit Metrik $ds^2 = (dx^2+dy^2)/y^2$ ist und $\Gamma \subset \mathrm{PSL}(2,\RR)$
eine kokompakte diskrete torsionsfreie Untergruppe. Dann ist $M$ eine kompakte
hyperbolische Fl\"ache vom Geschlecht $g \ge 2$. Der Laplace--Beltrami-Operator
$\Delop$ agiert auf $L^2(M)$ als ein nicht-negativer selbstadjungierter Operator mit
rein diskretem Spektrum
\begin{equation}
0 = \lambda_0 < \lambda_1 \le \lambda_2 \le \dots,\qquad \lambda_j \to \infty.
\end{equation}
Indem wir $\lambda_j = 1/4 + r_j^2$ mit $r_j \in \RR \cup i[-1/2, 1/2]$ schreiben, kodieren die
spektralen Parameter $r_j$ das Quantenspektrum.

\subsection{Primitive geschlossene Geod\"aten}

Die Konjugationsklassen hyperbolischer Elemente in $\Gamma$ korrespondieren bijektiv
mit (unorientierten) geschlossenen Geod\"aten auf $M$. Eine Geod\"ate ist \emph{primitiv}, wenn sie
keine positive Iteration einer anderen ist. Bezeichne mit $\mathcal{P}_\Gamma$ die Menge
der primitiven geschlossenen Geod\"aten, und f\"ur $\gamma \in \mathcal{P}_\Gamma$ mit
$\ell_\gamma$ ihre L\"ange.

Der Primgeod\"atensatz \cite{Huber1959, Sarnak1981} besagt
\begin{equation}
\label{eq:PGT}
\pi_\Gamma(T) := \#\{\gamma \in \mathcal{P}_\Gamma : \ell_\gamma \le T\} \sim \frac{e^T}{T},\qquad T \to \infty.
\end{equation}
Dies ist das geod\"atische Analogon zum PNT. Entscheidend ist: \emph{die Dichte ist
exponentiell, nicht polynomiell}: $\pi_\Gamma(T)$ w\"achst wie $e^T/T$, w\"ahrend
$\pi(x)$ wie $x/\log x = (e^{\log x})/\log x$ w\"achst. Mit anderen Worten, das
``Primzahlen-L\"angenspektrum'' von Geod\"aten, auf der $\log$-Achse betrachtet, hat die Dichte
$e^T/T$ (in der Variablen $T = \ell_\gamma$), was vergleichbar ist mit der Dichte von
Primzahlen auf der $\log$-Achse (in der Variablen $u = \log p$) erst nach der
Variablensubstitution.

\subsection{Selbergsche Zetafunktion}

\begin{definition}
\label{def:Selberg-zeta}
Die Selbergsche Zetafunktion ist
\begin{equation}
\label{eq:Sel-def}
\Sel(s) = \prod_{\gamma \in \mathcal{P}_\Gamma}\prod_{k=0}^\infty \bigl(1 - e^{-(s+k)\ell_\gamma}\bigr),\qquad \mathrm{Re}(s) > 1.
\end{equation}
\end{definition}

$\Sel(s)$ setzt sich meromorph nach $\CC$ fort, mit:

\begin{itemize}
\item \emph{Trivialen Nullstellen} bei $s = -k$, $k \in \NN_0$ (und weiteren trivialen Nullstellen,
die mit der Topologie assoziiert sind; siehe \cite{Hejhal1976}).
\item \emph{Nicht-trivialen Nullstellen} bei $s = 1/2 \pm ir_j$ f\"ur jeden
Laplace-Eigenwert $\lambda_j = 1/4 + r_j^2 > 0$.
\item Einer Nullstelle bei $s = 1$ (von $\lambda_0 = 0$).
\end{itemize}

\begin{theorem}[Selbergsche spektrale Nullstellenbeschreibung und strenges Linienkriterium, \cite{Selberg1956}]
\label{thm:Sel-RH}
Die nicht-trivialen spektralen Nullstellen von $\Sel(s)$, die zu
Laplace-Eigenwerten $\lambda_j=1/4+r_j^2$ geh\"oren, liegen bei
$s=1/2\pm ir_j$. Daher liefern Eigenwerte $\lambda_j\ge 1/4$ Nullstellen auf
der Geraden $\mathrm{Re}(s)=1/2$, w\"ahrend Ausnahmeeigenwerte
$0<\lambda_j<1/4$ reelle Nullstellen $s=1/2\pm\delta_j$ im kritischen Streifen
erzeugen. Insbesondere gilt die strenge kritische-Linien-Aussage unter der
Spektrall\"uckenbedingung $\lambda_1(\Delop)\ge 1/4$.
\end{theorem}

Dies folgt aus der Nicht-Negativit\"at von $\Delop$: $\lambda_j \ge 0$ impliziert
$r_j \in \RR \cup i[-1/2, 1/2]$.  Nicht-negative reelle $r_j$ (d.h.\ $\lambda_j \ge
1/4$) ergeben Nullstellen auf der kritischen Linie $\mathrm{Re}(s)=1/2$.  Au\ss{}ergewöhnliche
Eigenwerte $\lambda_j \in (0,1/4)$ produzieren jedoch $r_j = i\delta_j$ mit
$\delta_j \in (0,1/2)$ und liefern dadurch Nullstellen bei $s = 1/2 \pm \delta_j$,
die \emph{nicht} auf der kritischen Linie liegen.  Selberg-RH im strengen
spektralen Sinne (alle nicht-trivialen spektralen Nullstellen im kritischen
Streifen auf $\mathrm{Re}(s)=1/2$) erfordert daher die
Spektrall\"uckenbedingung $\lambda_1(\Delop) \ge 1/4$.  Dies ist Selbergs
Eigenwert-Vermutung \cite{Selberg1965}: sie ist in der schw\"acheren Form
$\lambda_1 \ge 3/16$ f\"ur Kongruenzuntergruppen bewiesen, bleibt aber als
$\lambda_1 \ge 1/4$ in voller Allgemeinheit offen.  Die v2.0-Analyse in
Abschnitt~\ref{sec:v2-to-selberg} wird unter dieser Hypothese ausgef\"uhrt.

\begin{remark}[Funktionalgleichung]
\label{rem:func-eq}
Die vervollst\"andigte Selbergsche Zetafunktion erf\"ullt die Funktionalgleichung
$\Sel(s) = \Sel(1-s) \cdot \Xi_\Gamma(s)$, wobei $\Xi_\Gamma(s)$ ein
expliziter meromorpher Faktor ist, der von $\mathrm{Area}(M)$ und dem Geschlecht von
$M$ abh\"angt \cite{Hejhal1976}.  Insbesondere treten nicht-triviale Nullstellen paarweise
$\rho$, $1-\rho$ auf, was ihre Symmetrie bzgl. $\mathrm{Re}(s)=\tfrac{1}{2}$ best\"atigt.
\end{remark}

\subsection{Selbergsche explizite Formel}

Die Selbergsche Spurformel \cite{Selberg1956, Hejhal1976} liefert das
geod\"atische Analogon der Weilschen expliziten Formel. F\"ur ein Paar
$(h, \hat h)$ ausreichend schnell abfallender Funktionen, die durch Fouriertransformation
verbunden sind, lautet die Spurformel
\begin{equation}
\label{eq:Selberg-trace}
\sum_{j=0}^\infty h(r_j) = \frac{\mathrm{Area}(M)}{4\pi}\int_\RR h(r)\,r\tanh(\pi r)\,dr
+ \sum_{\gamma \in \mathcal{P}_\Gamma}\sum_{m=1}^\infty \frac{\ell_\gamma\,\hat h(m\ell_\gamma)}{2\sinh(m\ell_\gamma/2)}.
\end{equation}
Die linke Seite ist die \emph{spektrale Seite} (eine Summe \"uber Eigenwerte von
$\Delop$); die rechte Seite ist die \emph{geometrische Seite} (ein Identit\"atsterm
plus eine Summe \"uber Geod\"aten). Dies spielt exakt die Rolle von
\begin{equation*}
\sum_\rho \Phi(\rho) = \text{(archimedischer Term)} - \sum_p\sum_{m\ge 1}\frac{\log p}{p^{m/2}}(\hat\Phi(m\log p) + \hat\Phi(-m\log p))
\end{equation*}
in der Weilschen expliziten Formel: Primzahlen $\leftrightarrow$ primitive Geod\"aten,
$\log p \leftrightarrow \ell_\gamma$, nicht-triviale Nullstellen $\rho \leftrightarrow$
spektrale Parameter $r_j$.

% =============================================================================
\section{Anwendung von v2.0 auf Selberg}
\label{sec:v2-to-selberg}
% =============================================================================

\subsection[Selberg-Geod\"aten-Verschiebungen]{Selberg-Geod\"aten-Verschiebungen und\\ Verschiebungs-Parit\"ats-Lemma}

Setze $L = \log\lambda$. Auf $L^2([-L, L])$ definieren wir den Selberg-Primzahl-Verschiebungs-Operator
\begin{equation}
\label{eq:PSel}
(P^\Gamma_\lambda f)(t) = \sum_{\gamma \in \mathcal{P}_\Gamma, \ell_\gamma \le 2L}\sum_{m=1}^{\lfloor 2L/\ell_\gamma\rfloor}
\frac{\ell_\gamma}{2\sinh(m\ell_\gamma/2)}\bigl[f(t-m\ell_\gamma) + f(t+m\ell_\gamma)\bigr]\mathbf{1}_{[-L,L]}.
\end{equation}
In der Kosinus-/Sinus-Basis von Abschnitt~\ref{sec:v2-recollect} definieren wir die
\emph{Selberg-Differenzmatrix}
\begin{equation}
\label{eq:DNSel}
D^\Gamma_{nm}(r) = \langle\varphi_n, (S_{rL}+S_{-rL})\varphi_m\rangle - \langle\psi_n, (S_{rL}+S_{-rL})\psi_m\rangle.
\end{equation}
Dies ist \emph{dieselbe Matrix} $D_N(r)$ wie im Riemann-Setting: sie h\"angt
nur von der normierten Verschiebung $r$ ab, nicht von der Quelle der Verschiebung (Primzahl-L\"ange
oder Geod\"aten-L\"ange). Folglich gilt:

\begin{proposition}
\label{prop:SPL-transfer}
Das Verschiebungs-Parit\"ats-Lemma (Theorem~\ref{thm:SPL}) gilt w\"ortlich im Selberg-Setting:
f\"ur jedes $N \ge 3$, jedes $r \in (0,2)$ gilt
$\lambda_{\min}(D^\Gamma_N(r)) < 0$.
\end{proposition}

\begin{proof}
$D^\Gamma_N(r) = D_N(r)$ als Matrizen. Das Verschiebungs-Parit\"ats-Lemma h\"angt nur von
der harmonischen Struktur der Kosinus- und Sinus-Basen auf $[-L, L]$ ab, und von der
Tatsache, dass $S_{rL}+S_{-rL}$ die Summe zweier Verschiebungen um den gleichen Betrag in
entgegengesetzte Richtungen ist. Der zahlentheoretische oder geometrische Ursprung von $r$ ist
unbedeutend.
\end{proof}

\begin{remark}
Dies ist die algebraische Universalit\"at. Im v2.0-Rahmen f\"ur die Riemannsche
Zetafunktion ist das Verschiebungs-Parit\"ats-Lemma das Fundament, auf dem alles ruht. Es
\"ubertr\"agt sich auf \emph{jedes} Setting, in dem Nullstellen durch eine explizite
Formel mit reellen, positiven L\"angen parametrisiert sind, die die Rolle von $\log p$ spielen. Dies
umfasst Dirichlet L-Funktionen, Dedekind Zetafunktionen, Selberg Zeta,
Ruelle Zeta und im Prinzip jede Hecke- oder automorphe L-Funktion.
\end{remark}

\subsection{Selberg-Weil quadratische Form}

Auf $L^2([-L, L])$ definieren wir die \emph{Selberg--Weilsche quadratische Form}
\begin{equation}
\label{eq:QW-Sel}
\QW^\Gamma_\lambda(f) = \frac{\mathrm{Area}(M)}{4\pi}K_\mathrm{arch}^\Gamma(f,f)
+ \sum_{\gamma \in \mathcal{P}_\Gamma, \ell_\gamma\le 2L}\sum_{m=1}^\infty
\frac{\ell_\gamma}{2\sinh(m\ell_\gamma/2)}\langle T_{m\ell_\gamma} f + T_{-m\ell_\gamma}f, f\rangle,
\end{equation}
\begin{remark}[Trunkierung ist exakt]
\label{rem:truncation-exact}
F\"ur $f \in L^2([-L,L])$ und $|u| > 2L$ verschwindet das innere Produkt der Verschiebung:
$\langle T_u f + T_{-u}f,\, f\rangle = 0$, da $f$ auf $[-L,L]$ getragen wird,
w\"ahrend $T_{\pm u}f$ einen Tr\"ager hat, der au\ss{}erhalb von $[-L,L]$ verschoben ist.  Folglich
tr\"agt jede geod\"atische Iteration mit $m\ell_\gamma > 2L$ \emph{exakt} null
zu $\QW^\Gamma_\lambda$ bei, ohne Approximationsfehler.  Die Summe in
\eqref{eq:QW-Sel} ist daher keine Approximation, sondern ein exakter Ausdruck
f\"ur die quadratische Form auf $L^2([-L,L])$.
\end{remark}

wobei $K_\mathrm{arch}^\Gamma$ das archimedische Analogon
\begin{equation}
\label{eq:arch-Sel}
K_\mathrm{arch}^\Gamma(f,f) = \int_0^\infty[\langle T_u f + T_{-u}f, f\rangle - 2\,c_\Gamma^{\mathrm{id}}(u)\|f\|^2]\cdot \kappa^\Gamma(u)\,du,
\end{equation}
ist, das durch die \"Ubersetzung des Identit\"atsterms in \eqref{eq:Selberg-trace} in
Verschiebungs-Operatoren erhalten wird.
\begin{remark}[Konvergenz und Renormierung bei $u=0$]
\label{rem:arch-renorm}
Gleichung~\eqref{eq:arch-Sel} ist der Ausdruck des Identit\"atsterms der
Selbergschen Spurformel im $u$-Raum angewandt auf $h(r) := |\hat f(r)|^2$; f\"ur
$f$ in einem dichten Unterraum (z.B.\ $C_c^\infty([-L,L])$) konvergieren
beide Seiten absolut.  Die punktweise Singularit\"at von $\kappa^\Gamma$ bei $u=0$
wird durch die von $c_\Gamma^{\mathrm{id}}$ ausgeglichen, und die Klammer ist der
regularisierte Identit\"atsbeitrag; vgl.\ \cite[Kap.~6]{Hejhal1976}.
\end{remark}
Der Koeffizient
\begin{equation}
\label{eq:cGamma-def}
c_\Gamma^{\mathrm{id}}(u) \;:=\; \frac{\mathrm{Area}(M)}{4\pi}
\int_0^\infty r\tanh(\pi r)\cos(ru)\,dr
\end{equation}
ist die Fourier-Kosinus-Transformierte der Plancherel-Ma\ss{}dichte
$r\tanh(\pi r)$, die den Beitrag der Identit\"atsklasse aus
\eqref{eq:Selberg-trace} kodiert; das Integral ist im Sinne einer temperierten Distribution zu
verstehen, da $r\tanh(\pi r)$ polynomiell
beschr\"ankt ist (vgl.\ \cite[Kap.~6]{Hejhal1976}).  Hierbei ist
\begin{equation}
\label{eq:kappa-Sel}
\kappa^\Gamma(u) := \frac{1}{4\pi}\,\frac{d}{du}\!\Big[\frac{1}{\sinh(u/2)}\Big],\qquad u > 0,
\end{equation}
der Plancherel-Kern der Selberg-Transformation \cite[Gl.~(1.40)]{Iwaniec2002}.
Man beachte: f\"ur alle $u > 0$ gilt $\kappa^\Gamma(u) < 0$, denn
\[
\frac{d}{du}\bigl[\sinh(u/2)^{-1}\bigr]
= -\frac{\cosh(u/2)}{2\sinh^2(u/2)} < 0.
\]
Dies ist konsistent mit der Positivit\"atsstruktur von $K_\mathrm{arch}^\Gamma$:
die Klammer in \eqref{eq:arch-Sel} ist negativ, wenn $c_\Gamma^{\mathrm{id}}(u)$ den
Identit\"atsoperator normiert, und das Produkt zweier negativer Gr\"o\ss{}en ist positiv.

Die quadratische Form $\QW^\Gamma_\lambda$ l\"asst eine Parit\"atszerlegung zu. Ihr
minimaler Eigenwert wird entweder durch eine gerade oder eine ungerade Eigenfunktion angenommen.

\begin{definition}[Selberg gerade Dominanz]
\label{def:Sel-ED}
Wir sagen, dass \emph{gerade Dominanz} (even dominance) f\"ur $\QW^\Gamma_\lambda$ beim Parameter
$\lambda$ gilt, wenn der minimale Eigenwert von $\QW^\Gamma_\lambda$ einfach ist und
von einer geraden Eigenfunktion angenommen wird.
\end{definition}

\begin{remark}[Spurformel als Br\"ucke zwischen spektralem und geometrischem Raum]
\label{rem:two-spaces}
Die Form $\QW^\Gamma_\lambda$ ist vollst\"andig im spektralen Parameterraum
$L^2([-L,L])$ definiert, wobei $T_u$ die trunkierte Verschiebung $f(\cdot + u)$
eingeschr\"ankt auf $[-L,L]$ bezeichnet.  Die Selbergsche Spurformel
\eqref{eq:Selberg-trace} koppelt diese Testfunktionsseite an die geometrische
Geod\"atensumme und an die Spektralsumme $\sum_j h(r_j)$.  Diese Kopplung ist
eine Gleichheit von spektralen Distributionen, nicht schon die Konstruktion
eines unit\"aren Intertwiners zwischen einer Geod\"atenfluss-Darstellung auf
$L^2(SM)$ und dem trunkierten Intervall-Shift auf $L^2([-L,L])$.  Alle
NE-B-Aussagen unten sind daher auf den sph\"arischen Spurformelkanal beschr\"ankt,
solange ein solcher expliziter Intertwiner nicht zus\"atzlich geliefert wird.
\end{remark}

\subsection{Frontier-Dominanz im exponentiellen Dichteregime}

Das Frontier-Band f\"ur Selberg ist die Menge von Geod\"aten mit
$\ell_\gamma / L$ nahe bei $1$, d.h.\ $\ell_\gamma \in [L-1, L]$. Nach dem
Primgeod\"atensatz \eqref{eq:PGT} gilt
\begin{equation}
\label{eq:frontier-count-Sel}
\#\{\gamma : \ell_\gamma \in [L-1, L]\} = \pi_\Gamma(L) - \pi_\Gamma(L-1)
\sim \frac{e^L}{L}(1 - e^{-1}) = \frac{\lambda}{\log\lambda}(1 - e^{-1}).
\end{equation}
Dies ist \emph{dieselbe asymptotische Gr\"o\ss{}enordnung} wie die Riemannsche Frontier-Anzahl
$\pi(\lambda) - \pi(\lambda/e)$, bis auf die Konstante $(1-e^{-1}) \approx 0.632$.

Das Frontier-Geod\"atengewicht ist f\"ur $\ell_\gamma \sim L$ gegeben durch
\begin{equation}
\label{eq:frontier-weight-Sel}
w^\Gamma_\gamma := \frac{\ell_\gamma}{2\sinh(\ell_\gamma/2)} \sim \frac{L}{e^{L/2}} = \log\lambda \cdot \lambda^{-1/2}.
\end{equation}
Vergleiche dies mit dem Riemannschen Frontier-Primzahlgewicht
$w_p^{\mathrm{Riem}} = \log p / p^{1/2} \sim L \cdot \lambda^{-1/2}$ f\"ur
$p \sim \lambda$. Die Gewichte stimmen asymptotisch in ihrer Gr\"o\ss{}enordnung \"uberein.

Durch Kombination von \eqref{eq:frontier-count-Sel} und der Gewichtsabsch\"atzung mit dem
Verschiebungs-Parit\"ats-Lemma erhalten wir das Selberg-Analogon der
Frontier-Dominanz-Absch\"atzung, sobald eine quantitative SPL-Magnitudenschranke
und eine Selberg-Off-Frontier-/archimedische Remainder-Schranke angenommen werden:
\begin{equation}
\label{eq:frontier-bound-Sel}
\Delta^\Gamma(\lambda) := \lambda_1^+(\QW^\Gamma_\lambda) - \lambda_1^-(\QW^\Gamma_\lambda)
\le -c_\Gamma\sqrt{\lambda} + o(\sqrt{\lambda}),
\end{equation}
f\"ur eine positive Konstante $c_\Gamma > 0$, die von der Fl\"ache $M$ abh\"angt (durch
$\mathrm{Area}(M)$ und die Selbergsche Zetafunktion nahe $s=1$).

\begin{lemma}[Bedingte Frontier-Dominanz-Konstante]
\label{lem:frontier-constant}
Angenommen, die Selberg-Frontier-Summanden erf\"ullen die quantitative
SPL-Magnitudenschranke aus \cite[Lem.~A.3]{Geiger2026partII} nach dem
Selberg-Gewichtstransfer und die Nicht-Frontier- plus archimedischen Terme sind
kleinerer Ordnung als $\sqrt{\lambda}$. Dann existiert f\"ur jede kompakte
hyperbolische Fl\"ache $M$ ein $\lambda_0 = \lambda_0(M) > 0$, sodass die
Frontier-Dominanz-Konstante in \eqref{eq:frontier-bound-Sel} Folgendes erf\"ullt:
\begin{equation}
\label{eq:cGamma-explicit}
c_\Gamma \;\ge\; \tfrac{1}{2}(1-e^{-1}) \;>\; 0
\qquad\text{f\"ur alle }\lambda \ge \lambda_0.
\end{equation}
\end{lemma}

\begin{remark}
Die normierte vorzeichenlose Frontier-Summe $\lambda^{-1/2}\sum_{\ell_\gamma\in[L-1,L]}
w^\Gamma_\gamma \to (1-e^{-1}) \approx 0.632$ f\"ur $\lambda\to\infty$.
Der Faktor $\tfrac{1}{2}$ in \eqref{eq:cGamma-explicit} ist konservativ: er verwendet
die quantitative untere Schranke aus \cite[Lem.~A.3]{Geiger2026partII}, die
eine vorzeichenbehaftete Magnitude $\ge\tfrac{1}{2}w^\Gamma_\gamma$ pro Frontier-Geod\"ate liefert.
Proposition~\ref{prop:SPL-transfer} liefert allein das Vorzeichen ($\lambda_{\min}<0$);
die Magnituden-Schranke und die globale Remainder-Kontrolle sind die bedingten
Inputs dieses Lemmas.
Unter der scharfen Schranke geht $c_\Gamma \to (1-e^{-1})$ asymptotisch.
\end{remark}

\begin{proof}[Beweisskizze]
Unter den genannten Magnituden- und Remainder-Annahmen enth\"alt das
Frontier-Band $[L-1,L]$
$(1-e^{-1})\,e^L/L\,(1+o(1)) = (1-e^{-1})\,\lambda/\log\lambda\,(1+o(1))$
primitive Geod\"aten nach \eqref{eq:frontier-count-Sel}. Jede tr\"agt das Gewicht
$w^\Gamma_\gamma = L/(2\sinh(L/2)) \sim \log\lambda\cdot\lambda^{-1/2}$ nach
\eqref{eq:frontier-weight-Sel} bei. Nach \cite[Lem.~A.3]{Geiger2026partII}
ist jeder Frontier-Summand f\"ur alle gro\ss{}en $\lambda$ negativ mit einer Gr\"o\ss{}e von mindestens
$\tfrac{1}{2}w^\Gamma_\gamma$.  Die Summation \"uber das Frontier-Band
und Division durch $\sqrt\lambda$ ergibt
\[
\sum_{\ell_\gamma\in[L-1,L]} w^\Gamma_\gamma \;\ge\;
(1-e^{-1})\,\frac{\lambda}{\log\lambda}\cdot
\frac{\log\lambda}{\sqrt\lambda}\cdot\tfrac{1}{2}\,(1+o(1))
= \tfrac{1}{2}(1-e^{-1})\sqrt\lambda\,(1+o(1)),
\]
was $c_\Gamma \ge (1-e^{-1})/2$ liefert.
\end{proof}

\begin{remark}
Der Faktor $(1-e^{-1}) \approx 0.632$ ist das geometrische Verh\"altnis der Frontier-Geod\"atendichte
(exponentielles Regime) zur Frontier-Primzahldichte
(polynomielles Regime): Selberg hat $(1-e^{-1})\lambda/\log\lambda$ Frontier-Geod\"aten,
wo Riemann $\lambda/\log\lambda$ Frontier-Primzahlen hat.
Dies erkl\"art den einzigen strukturellen Unterschied zwischen den beiden Frontier-Absch\"atzungen.
\end{remark}

\subsection{Bedingte Reduktion \"uber v2.0}

\begin{theorem}[v2.0 Selberg-RH, bedingt]
\label{thm:v2-Sel-RH}
Sei $M = \Gamma\backslash\HH$ eine kompakte hyperbolische Fl\"ache mit
$\lambda_1(\Delop) \ge 1/4$.  Angenommen:
\begin{enumerate}[label=(C2\alph*)]
  \item Die Selberg-Quadratform $\QW^\Gamma_\lambda$ besitzt eine nach unten
    beschr\"ankte selbstadjungierte Intervall-Realisierung.
  \item F\"ur alle $\lambda \ge \lambda_0(\Gamma)$ ist der minimale Eigenwert
    von $\QW^\Gamma_\lambda$ einfach und wird von einer geraden Eigenfunktion
    angenommen.
  \item Die Fourier-/Selberg-Transformierte dieses Minimizers ist mit dem
    abgeschlossenen Selberg-Spektralfaktor identifiziert, dessen Nullstellen die
    nicht-trivialen spektralen Nullstellen von $\Sel(s)$ parametrisieren.
\end{enumerate}
Dann liegen alle nicht-trivialen spektralen Nullstellen von $\Sel(s)$ im
kritischen Streifen auf $\mathrm{Re}(s) = 1/2$. Ohne (C2c) liefert Connes'
Mechanismus nur eine Aussage \"uber die Nullstellen der Minimizer-Transformierten,
nicht \"uber die spektralen Nullstellen von $\Sel$.
\end{theorem}

Zur Leserführung wird das C2-Gate getrennt vom Theorem dokumentiert.
Tabelle~\ref{tab:c2-minimizer-ledger} unterscheidet, welche Bestandteile in
diesem Papier bereits verwendet werden und welche als Companion-Audit offen
bleiben, statt eine zusätzliche Theorem-Schließung zu behaupten. Die
Bedingung ohne Ausnahmeeigenwerte wird als separates SG/NoExc-Gate geführt;
sie ist nicht Teil von C2b.

\begin{table}[htp]
\centering
\small
\setlength{\tabcolsep}{3pt}
{\renewcommand{\arraystretch}{1.16}
\begin{tabularx}{\textwidth}{@{}L{0.18\textwidth}YYY@{}}
\toprule
Gate & Dokumentiertes Objekt & Aktueller Einsatz im Papier & Fehler- oder Revisionssignal \\
\midrule
SPH: sphärischer Kanal
  & Selberg-Transformation und radiale \(K\)-bi-invariante Faltung über
    \(\widehat{k}(r_j)\), mit Casimir-/Laplace-Kommutation.
  & Positivkontroll-Kontext für das Fehlen von NE-B; selbst kein C2-Gate und
    keine Identifikation individueller Geodätenfluss-Pullbacks mit
    Intervallverschiebungen.
  & Ein nicht-radiales oder volles \(SM\)-Shift-Modell wird wie das trunkierte
    \(L^2([-L,L])\)-Modell behandelt. \\
C2a: Intervall-Realisierung
  & Nach unten beschränkte selbstadjungierte Intervall-Realisierung von
    \(\QW^\Gamma_\lambda\).
  & Explizite v2.0-Hypothese für die Anwendung des Connes-Mechanismus.
  & Das Intervallmodell wird nur per Analogie oder aus dem sphärischen Kanal
    behauptet, ohne eine angegebene selbstadjungierte Realisierung. \\
C2b: Paritäts-/Minimizer-Gate
  & Einfacher minimaler Eigenwert und gerader Minimizer der
    Intervall-Realisierung.
  & Explizite v2.0-Hypothese; durch den Selberg-/Casimir-Kanal gestützt, aber
    hier nicht aus der Spektrallücke abgeleitet.
  & Das Minimum ist entartet, ungerade oder erst nach Beobachtung des
    Ziel-Selberg-Faktors ausgewählt. \\
SG/NoExc: Spektrallinien-Input
  & \(\lambda_1(\Delop)\ge 1/4\) und Abwesenheit von Ausnahmeeigenwerten.
  & Explizite Hypothese für die strenge kritische-Linien-Form; die spektrale
    Nullstellenbeschreibung bleibt klassisch.
  & Ein Ausnahmeeigenwert \(0<\lambda_j<1/4\) erzeugt reale Nullstellen
    \(1/2\pm\delta_j\). \\
C2c: Minimizer-/Fenster-Identifikation
  & Zielfenster \(I\), Zielquelle, Zentrum \(\nu\), äußerer Gap, Residuum, Masse
    außerhalb des Fensters und Negativkontrolle.
  & Operatives Ledger inzwischen mit realen Spektraldaten kalibriert
    (\Cref{sec:c2c-real-ledger}); die Connes-Minimizer-Identifikation selbst
    bleibt offen, C2c bleibt Companion-/Audit-Gate.
  & Richtiger Rayleigh-Schwerpunkt, aber falscher Tail oder Masse außerhalb des
    Fensters (\texttt{same\_rayleigh\_}\allowbreak\texttt{wrong\_tail}). \\
Entscheidung
  & \texttt{pass}, \texttt{ambiguous} oder \texttt{fail} je
    Flächen-/Fensterkandidat.
  & Positivkontroll-Audit, kein Beweis-Upgrade; lokalisiert den Rest für
    JST-/Companion-Arbeit.
  & Fehlende Tabelle, Zielwahl aus demselben Lauf oder fehlende Negativkontrollen. \\
\bottomrule
\end{tabularx}}
\caption{C2-/Minimizer-Identifikationsledger für die Selberg-Positivkontrolle.}
\label{tab:c2-minimizer-ledger}
\end{table}

\begin{proof}[Beweisskizze]
Die Struktur des Arguments ist identisch zum v2.0-Beweis f\"ur Riemann.

\smallskip\noindent
\textbf{Schritt 1 (punktweise Bevorzugung von geraden Funktionen, einschlie\ss{}lich Iterationen).} Nach
Proposition~\ref{prop:SPL-transfer} gilt f\"ur jedes $\gamma \in \mathcal{P}_\Gamma$
und jedes $m \ge 1$ mit $m\ell_\gamma \le 2L$, dass die Differenzmatrix
$D^\Gamma_N(m\ell_\gamma / L)$ $\lambda_{\min} < 0$ erf\"ullt.  Hier entspricht $m > 1$
der $m$-fachen Iteration von $\gamma$; das SPL ist auf Iterationen genauso
anwendbar wie auf primitive Geod\"aten, da die Verschiebungsl\"ange $m\ell_\gamma$ in das SPL als skalarer
Parameter $r = m\ell_\gamma/L$ eingeht.
Nach Bemerkung~\ref{rem:truncation-exact} tragen Iterationen mit $m\ell_\gamma > 2L$
exakt null bei und m\"ussen nicht separat behandelt werden.  Somit tr\"agt jeder einzelne
Geod\"aten-Summand (primitiv oder iteriert), gewichtet mit dem positiven Koeffizienten
$\ell_\gamma/(2\sinh(m\ell_\gamma/2))$, negativ zur Gerade-minus-Ungerade-Differenz bei.

\smallskip\noindent
\textbf{Schritt 2 (Frontier-Dominanz).} Unter den Annahmen von
Lemma~\ref{lem:frontier-constant} ist der aggregierte Beitrag von
Frontier-Geod\"aten zu $\Delta^\Gamma(\lambda)$ negativ von der Ordnung
$-c_\Gamma \sqrt{\lambda}$, w\"ahrend der kombinierte Beitrag der
Nicht-Frontier-Geod\"aten plus des archimedischen Terms kleinerer Ordnung ist.

\smallskip\noindent
\textbf{Schritt 3 (C2b: Paritäts-/Minimizer-Gate).} Die Einfachheit des minimalen
Eigenwerts und die Geradheit des Minimizers sind nicht automatisch eine Folge
von generischer Laplace-Eigenwert-Einfachheit. Sie sind genau das Gate (C2b).
Die Positivit\"at des hyperbolischen Kerns und der sph\"arische Casimir-Kanal
liefern starke strukturelle Evidenz, ersetzen aber nicht die explizite
Minimizer-Identifikation. C2b ist nicht das separate SG/NoExc-Spektrallücken-Gate.

\smallskip\noindent
\textbf{Schritt 4 (Connes' Reduktion).} Theorem~\ref{thm:connes} ist der
\emph{algebraische Mechanismus}: Gegeben sind (i) ein nach unten beschr\"ankter
selbstadjungierter Operator auf einem symmetrischen Intervall und (ii) ein
einfaches Minimum mit gerader Eigenfunktion, dann liegen alle Nullstellen der
Fouriertransformierten des minimalen Eigenvektors auf der reellen Achse.  Die
Verifikation von (i) und (ii) ist der $C2^\Gamma$-Teilschritt, w\"ahrend die
Identifikation des resultierenden Minimizer-Transforms mit dem abgeschlossenen
Selberg-Faktor die Bedingung (C2c) ist.  (In der
Drei-Komponenten-Zerlegung $\mathrm{RH}_X = \mathrm{GBC}_X + C2_X + \mathrm{Hurwitz}_X$
aus der Begleitarbeit \cite{GeigerMeta2026}, in Vorbereitung, entspricht dies
dem Lie-Ursprungs-Zweig plus dem expliziten Minimizer-Identifikationsgate.)
Im Riemann-Fall ist insbesondere Bestandteil~(ii) nichttrivial; genau dies ist
der Gehalt von v2.1 Proposition~A6 \cite{Geiger2026partII}.  Wendet man die
Connes-Reduktion mit $\QW^\Gamma_\lambda$ anstelle von $\QW_\lambda$ an, so
sind die Nullstellen von $\widehat{\eta}_{\lambda,\Gamma}$ reell.  Die
zus\"atzliche Bedingung (C2c) transferiert diese Realnullstellen-Aussage auf den
abgeschlossenen Selberg-Faktor. Zusammen mit $\lambda_1(\Delop)\ge 1/4$ ergibt
dies die strenge Aussage für die nicht-trivialen spektralen Selberg-Nullstellen
im kritischen Streifen unter den expliziten Hypothesen des Theorems.
\end{proof}

\begin{remark}[Bedingtheit und Konsistenz]
Theorem~\ref{thm:v2-Sel-RH} ist eine bedingte Reduktion: \emph{gegeben} die
C2-Gates und die separate Bedingung ohne Ausnahmeeigenwerte, folgt die strenge
spektrale Nullstellenaussage durch das v2.0-Argument.  Da die klassische
Selberg-Theorie bereits die spektrale Nullstellenbeschreibung liefert und die
strenge Linienform an die L\"ucke $\lambda_1(\Delop)\ge 1/4$ bindet, verifiziert
sie die Schlussfolgerung auf der richtigen Hypothesenebene unabh\"angig.  Der
Wert der bedingten Reduktion liegt in der \emph{Belastungsprobe der Reichweite
der v2.0-Methode}: sie ist mit einem Fall kompatibel, in dem ein
Hilbert--P\'olya-/Casimir-Kanal bereits existiert. Diese \"Ubereinstimmung ist
ein Konsistenzbeleg, kein zus\"atzlicher unbedingter Selberg-RH-Beweis.
\end{remark}

\begin{remark}[Gerade Dominanz als strukturelle Erwartung]
\label{rem:even-dominance-structural}
Die Annahme der geraden Dominanz in Theorem~\ref{thm:v2-Sel-RH} ist eine
\emph{strukturelle Erwartung}, aber keine hier bewiesene automatische
Konsequenz aus der Laplace-Spektraltheorie.  Sie ist im Theorem als Gate (C2b)
ausdr\"ucklich angenommen.  Ein vollst\"andig strenger Beweis w\"urde dieselben
Frontier-Dominanz- und Remainder-Absch\"atzungen erfordern wie im
Riemann-Teil~II \cite{Geiger2026partII}, angepasst an die exponentielle
Geod\"aten-Dichte, plus die Minimizer-Identifikation (C2c).  Der sph\"arische
Casimir-Kanal st\"utzt die Konsistenz, ersetzt aber nicht diese Gates.
\end{remark}

\subsection{C2c-Fensterledger mit realen Spektraldaten}
\label{sec:c2c-real-ledger}

Das C2c-Gate aus Tabelle~\ref{tab:c2-minimizer-ledger} verlangt ein
operatives Fensterzertifikat. Für einen Kandidatenvektor $\psi$, ein
Spektralfenster $I$ mit Zentrum $\nu$ und den relevanten selbstadjungierten
Operator $D$ (hier $D=\Delop$ bzw.\ dessen kuspidale Restriktion) liefert
der Spektralsatz
\begin{equation}
\label{eq:window-certificate}
\|(\mathrm{Id}-P_I)\psi\| \;\le\;
\frac{\|(D-\nu)\psi\|}{\operatorname{dist}\bigl(\nu,\,\sigma(D)\setminus I\bigr)},
\qquad P_I = \mathbf{1}_I(D).
\end{equation}
Dieser Abschnitt berichtet eine Kalibrierung dieses Zertifikats auf
\emph{realen, unabhängig berechneten} Spektren; sie ersetzt das interne
Toy-Spektrum des früheren Projekt-Audits. Verwendet werden zwei
Positivkontrollflächen:
\begin{itemize}
\item die \emph{Bolza-Fläche} (kompakt, Genus $2$): die ersten zehn
positiven Laplace-Eigenwerte mit Multiplizitäten aus der zertifizierten
Berechnung von Strohmaier--Uski \cite{StrohmaierUski2013}. Kompaktheit macht
dies zum Lemma-konformen Hauptfall; das Verfahren aus
\cite{StrohmaierUski2013} findet \emph{alle} Eigenwerte unterhalb des
Cutoffs, die endliche Liste ist also unterhalb von
$\lambda_{\max}\approx 30{,}83$ vollständig;
\item die \emph{Modulfläche} $\mathrm{PSL}(2,\ZZ)\backslash\HH$
(Level $1$): die ersten acht kuspidalen Spektralparameter $r$
($\lambda = 1/4 + r^2$) aus den rigorosen Maass-Berechnungen von
Seymour-Howell \cite{SeymourHowell2022} und der in \cite{LowryDuda2025}
beschriebenen rigorosen Datenbank. Diese Fläche ist \emph{nicht} kompakt:
alle zugehörigen Zeilen beziehen sich ausschließlich auf die
selbstadjungierte Restriktion $\Delta_{\mathrm{cusp}}$ auf
$L^2_{\mathrm{cusp}}$; das kontinuierliche Eisenstein-Spektrum
$[1/4,\infty)$ des vollen Laplace-Operators liegt außerhalb des
Ledger-Scopes.
\end{itemize}

Das Zertifikat \eqref{eq:window-certificate} wird als
\emph{Finite-Spectrum-Proxy mit explizitem Tail-Budget} ausgewertet. Für
Multiplikator-Kandidaten $h(D)\psi_0$ wird die Spektralmasse oberhalb des
Daten-Cutoffs $\Lambda$ durch numerisch integrierte Tail-Budgets
$T_0 \ge \sum_{\lambda>\Lambda}\mathrm{mult}\cdot|h(\lambda)|^2$ und
$T_2 \ge \sum_{\lambda>\Lambda}\mathrm{mult}\cdot|h(\lambda)|^2(\lambda-\nu)^2$
beschränkt, unter einer dokumentierten Zähldichte-Annahme oberhalb des
Cutoffs (zweifacher Weyl-Hauptterm $\mathrm{Area}/4\pi$: $\rho=2$ für Bolza,
$\rho=1/6$ für den kuspidalen Modulflächenfall). Die berichteten
Residuum/Gap- und Off-Window-Spalten (beides Normgrößen, passend zur
Normform des Zertifikats) sind mit diesen Budgets worst-case-adjustiert, und ein Bracketing-Check verifiziert
$\Lambda-\nu > \operatorname{dist}(\nu,\sigma(D)\setminus I)$ für jedes
Fenster, sodass Eigenwerte oberhalb des Cutoffs den berichteten Gap nicht
verkleinern können. Die Audit-Schwellen (\texttt{pass} $\le 0{,}25$,
\texttt{ambiguous} $\le 0{,}5$) sind eine feste Konvention des Audits, nicht
Teil der Spektralsatz-Schranke.

Die Kandidaten sind in vier Kontrollklassen gruppiert:
\emph{Positivkontrollen} (Fenster-Eigenvektoren und getunte
Fenster-Multiplikatoren $h(D)$), \emph{radiale Fehlkandidaten} (echte
Funktionen von $D$ --- also exakt kommutierend ---, aber falsch zentriert
oder schlecht getunt), \emph{mathematische Negativkontrollen} (adversarielle
Spektralverteilungen, einschließlich der Pflichtklasse
\texttt{same\_rayleigh\_}\allowbreak\texttt{wrong\_tail}) und
\emph{Bookkeeping-Ablationen} (Strukturbrecher wie permutierte
Spektralbuchhaltung, die die Implementierung testen, nicht die
mathematische Kernaussage). Für Kandidaten der Form $h(D)$ verschwindet der
Kommutator-Leakage-Term identisch per Funktionalkalkül; das ist eine
\emph{Klassen}aussage, keine geprüfte Eigenschaft eines externen
Minimizer-Kandidaten.

\begin{table}[htp]
\centering
\footnotesize
\setlength{\tabcolsep}{4pt}
{\renewcommand{\arraystretch}{1.12}
\begin{tabularx}{\textwidth}{@{}L{2.5cm}L{4.1cm}L{1.9cm}rrl@{}}
\toprule
\textbf{Fenster} & \textbf{Kandidat} & \textbf{Klasse} &
\textbf{Res./Gap} & \textbf{Off-Win.} & \textbf{Entscheidung} \\
\midrule
Bolza $\lambda_1$ (Mult.\ $3$)
  & Fenster-Eigenraum & positiv & 0.000 & 0.000 & pass \\
  & getunter Gauss-Multiplikator & positiv & 0.002 & 0.002 & pass \\
  & ungetunter Heat-Kernel $e^{-\Delop}$ & Fehlkand. & 2.533 & 0.999 & fail \\
  & Heat-Kernel auf $L^2_0$ & positiv & 0.248 & 0.246 & pass \\
  & breiter Gauss ($\sigma{=}5$) & Fehlkand. & 1.123 & 0.758 & fail \\
  & dezentrierter Bandpass (bei $10$) & Fehlkand. & 2.912 & 1.000 & fail \\
  & same-Rayleigh wrong-tail & negativ & 1.592 & 1.000 & fail \\
  & falsches Zentrum & negativ & 1.000 & 1.000 & fail \\
  & gebrochene radiale Kommutanz & Ablation & 7.188 & 1.000 & fail \\
  & permutiertes Spektrum & Ablation & 16.003 & 1.000 & fail \\
\midrule
Bolza $\lambda_8$ (isoliert)
  & Fenster-Eigenraum & positiv & 0.000 & 0.000 & pass \\
  & getunter Gauss-Multiplikator & positiv & 0.003 & 0.003 & pass \\
  & same-Rayleigh wrong-tail & negativ & 1.398 & 1.000 & fail \\
\midrule
Modulfläche $\lambda_2$ (kuspidal)
  & Fenster-Eigenvektor & positiv & 0.000 & 0.000 & pass \\
  & getunter Gauss-Multiplikator & positiv & 0.000 & 0.000 & pass \\
  & same-Rayleigh wrong-tail & negativ & 1.172 & 1.000 & fail \\
\bottomrule
\end{tabularx}}
\caption{C2c-Fensterledger auf realen Spektren (Rev.~2; Residuum/Gap- und
Off-Window-Normwerte tail-adjustiert). Fenster: Bolza $\lambda_1$
(Multiplizität $3$), Bolza $\lambda_8$ (isoliert), Modulfläche $\lambda_2$
(kuspidal; $\Delta_{\mathrm{cusp}}$-Scope).}
\label{tab:c2c-real-ledger}
\end{table}

Das Ergebnis ist genau das Kalibrierungsmuster, das die C2c-Zeile von
Tabelle~\ref{tab:c2-minimizer-ledger} verlangt: Alle Positivkontrollen
bestehen auch nach Tail-Adjustierung (die Budgets der getunten Gauss-Fenster
sind numerisch vernachlässigbar; der auf $L^2_0$ eingeschränkte Heat-Kernel
--- zulässig, da $L^2_0$ $\Delop$-invariant ist --- bleibt ein
Grenzfall-Pass bei Ratio $0{,}248$); der \emph{ungetunte} Heat-Kernel
$e^{-\Delop}$ fällt durch, weil er auf den konstanten Grundzustand
$\lambda_0=0$ statt auf das Zielfenster konzentriert (eine korrekte
Diagnose, kein Fehlalarm); die beiden echten, aber falsch getunten
$h(D)$-Kandidaten scheitern \emph{innerhalb} der kommutierenden Klasse ---
Kommutation allein rettet also kein falsch gebautes Fenster; und die
Pflicht-Negativklasse
\texttt{same\_rayleigh\_}\allowbreak\texttt{wrong\_tail} scheitert auf
realen Daten trotz exakt getroffenen Rayleigh-Schwerpunkts. Das
Ledger-Skript und die vollständigen Ergebnisdateien (einschließlich der
$T_0/T_2$-Tail-Spalten und des Bracketing-Checks) sind im Projekt-Repository
verfügbar.\footnote{\url{https://github.com/research-line/functional-stability-theory/tree/main/masters/selberg}}

\begin{remark}[Kein Claim-Upgrade]
\label{rem:ledger-no-upgrade}
Dieses Ledger operationalisiert und kalibriert das C2c-\emph{Zertifikat} auf
realen Spektren; es schließt nicht das C2c-\emph{Gate}. Offen bleibt die
Kern-Identifikation: den Connes-Minimizer $\eta_{\lambda,\Gamma}$ selbst als
Vektor zu konstruieren, nachzuweisen, dass er in der $h(\Delop)$-Klasse
liegt (sodass die Klassenaussage zum Kommutator-Leakage auf ihn zutrifft),
und ihn durch dieses Ledger zu schicken. Bis dahin behält C2c
Companion-/Audit-Status, die Zähldichte-Schranke oberhalb des Cutoffs bleibt
eine dokumentierte Annahme (zweifacher Weyl-Hauptterm) und keine zitierte
Restterm-Schranke, und keine Schlussfolgerung dieses Papiers wird
aufgewertet.
\end{remark}

\subsection{NE-A f\"ur Selberg}

Wir definieren den Selberg-Primzahl-Verschiebungs-Multiplikator
\begin{equation}
\label{eq:Mxi-Sel}
M^\Gamma(\xi) = -2\,\mathrm{Re}\Big[\frac{\Sel'}{\Sel}\Big(\tfrac{1}{2}+i\xi\Big)\Big].
\end{equation}

\begin{proposition}[NE-A f\"ur Selberg]
\label{prop:NEA-Sel}
Die Menge $\{\xi \in \RR : M^\Gamma(\xi) < 0\}$ hat ein positives Lebesgue-Ma\ss{}.
\end{proposition}

\begin{proof}[Beweisskizze]
Aus der Hadamard-Produkt-Darstellung der Selbergschen Zetafunktion auf einer
kompakten hyperbolischen Fl\"ache (\cite[Kap.~10]{Iwaniec2002}) ergibt sich
\begin{equation}
\label{eq:hadamard-Sel}
\begin{aligned}
\frac{\Sel'}{\Sel}(s)
&= b_\Gamma + \sum_{\rho}\Big[\frac{1}{s-\rho}+\frac{1}{\rho}\Big] \\
&\quad + (\text{Beitr\"age von trivialen Nullstellen und dem topologischen Faktor}),
\end{aligned}
\end{equation}
wobei $\rho$ \"uber die nicht-trivialen Nullstellen l\"auft (auf $\mathrm{Re}(s)=1/2$ nach
Theorem~\ref{thm:Sel-RH}) und $b_\Gamma \in \RR$ eine Konstante ist, die von $\Gamma$ abh\"angt.

Setzt man $s = 1/2 + i\xi$ und nimmt den Realteil, so liefert jede nicht-triviale Nullstelle
$\rho_j = 1/2 + ir_j$ einen Beitrag
\begin{equation}
\label{eq:NEA-re}
\mathrm{Re}\Big[\frac{1}{i(\xi-r_j)}+\frac{1}{1/2+ir_j}\Big]
  = \frac{1/2}{1/4+r_j^2} > 0.
\end{equation}
Somit erzeugen die nicht-trivialen Nullstellen einen
\emph{positiven konstanten} Beitrag; sie verursachen keine Vorzeichenoszillationen
in $M^\Gamma(\xi)$.

Die Vorzeichenoszillationen entstehen durch die trivialen Nullstellen (bei $s = -k$, $k \ge 0$)
und den topologischen Faktor, deren Beitr\"age zu $\mathrm{Re}[\Sel'/\Sel]$
die Form $\sum_{k\ge 0} c_k/[(\xi^2+(k+1/2)^2)]$ mit $c_k \in \RR$
mit wechselndem Vorzeichen haben. Durch Kombination mit dem archimedischen Identit\"atsterm der
expliziten Selberg-Spurformel \eqref{eq:Selberg-trace} findet man durch direkte
Anwendung mit einer geeigneten Testfunktion (lokalisierend in $\xi$), dass
$M^\Gamma(\xi)$ sowohl strikt positive als auch strikt negative Werte auf
Mengen positiven Lebesgue-Ma\ss{}es annimmt. Das Argument ist das direkte Analogon zu
\cite[Lem.~A.3]{Geiger2026partII}; die einzige Ersetzung ist $\tanh(\pi r)$
anstelle von $\Gamma'/\Gamma(1/4+ir/2)$ im archimedischen Kern.
\end{proof}

NE-A l\"asst sich also \"ubertragen, aber mit der in der Beweisskizze sichtbaren
Rollenverteilung: Im Selberg-Fall stammen die Vorzeichenwechsel aus den
trivialen, topologischen und Identit\"atsbeitr\"agen der Spurformel, nicht aus den
kritischen-Linien-Nullstellen selbst.  Die generische Aussage ist daher
vorsichtiger: explizite Formeln enthalten typischerweise oszillatorische
logarithmische Ableitungsbeitr\"age, deren Quelle je nach Zeta-Familie
unterschiedlich ist.

% =============================================================================
\section[NE-B scheitert f\"ur Selberg: Der sph\"arische Casimir-Kanal]{NE-B scheitert f\"ur Selberg: Der sph\"arische Casimir-Kanal}
\label{sec:ne-b-fails}
% =============================================================================

Hier kommen wir zum wesentlichen Unterschied zwischen Selberg und Riemann.

\subsection{Aussage}

\begin{proposition}[Casimir-Kommutation im sph\"arischen Selberg-Kanal]
\label{thm:NEB-fail}
Seien $G=\mathrm{PSL}(2,\RR)$, $K=\mathrm{SO}(2)$ und
$M=\Gamma\backslash G/K$ eine kompakte hyperbolische Fl\"ache.  Sei $\Delop$
der Laplace--Beltrami-Operator, identifiziert mit dem Casimir-Element auf dem
sph\"arischen Unterraum.  F\"ur jeden radialen bi-$K$-invarianten Faltungskern
$k\in C_c^\infty(K\backslash G/K)$ kommutiert der Faltungsoperator
$R_k$ mit $\Delop$:
\begin{equation}
\label{eq:NEB-commut}
[\Delop, R_k] = 0.
\end{equation}
Folglich besitzt der sph\"arische Spurformelkanal einen nichttrivialen
universellen Kommutanten. Das relevante Selberg-Analogon von NE-B
(Theorem~\ref{thm:NEB-riemann}) ist in diesem pr\"azisen Kanal nicht vorhanden.
\end{proposition}

\subsection[Der sph\"arische Faltungsoperator auf L2(M)]{Der sph\"arische Faltungsoperator auf \texorpdfstring{$L^2(M)$}{L2(M)}}

Die Objekte, die v2.0 manipuliert, sind Verschiebungen in der \emph{logarithmischen Variablen}
$t \in [-L, L]$. Im Riemann-Setting entspricht die Variable $t$
$\log x$ im Idelklassen-Raum (in Connes' adelischem Rahmen). Im
Selberg-Setting wird die analoge L\"angenvariable im sph\"arischen
Spurformelkanal durch radiale Faltungskerne umgesetzt.

F\"ur einen kompakt getragenen $K$-biinvarianten Kern
$k\in C_c^\infty(K\backslash G/K)$ sei
\[
\begin{aligned}
(R_k f)(x)&=\int_G k(g)f(xg)\,dg,\\
f&\in L^2(\Gamma\backslash G/K).
\end{aligned}
\]
auf dem automorphen Quotienten definiert, wann immer das Integral existiert.
Diese radialen Operatoren sind die Spektralseite der Selberg-Spurformel: Die
Transformierte $\widehat{k}(r)$ diagonalisiert $R_k$ auf
Laplace-Eigenfunktionen:
\[
\Delop\psi_j=(1/4+r_j^2)\psi_j,
\qquad
R_k\psi_j=\widehat{k}(r_j)\psi_j.
\]

Die Spurformel schreibt die geod\"atische L\"angensumme in genau diesem
sph\"arischen Testfunktionskanal.  Das ist die richtige Selberg-Analogie zum
Kommutator-Test: nicht individuelle Pullbacks des geod\"atischen Flusses auf
$SM$, sondern die radialen Faltungsoperatoren, deren Selberg-Transform die
Spektralparameter $r_j$ sieht.

\subsection{Beweis von Proposition~\ref{thm:NEB-fail}}

\begin{proof}
Wir beweisen \eqref{eq:NEB-commut}, indem wir sowohl $\Delop$ als auch $R_k$
innerhalb der sph\"arischen Faltungsalgebra auf $\Gamma\backslash G/K$
identifizieren.

\smallskip\noindent
\textbf{Schritt 1 (Casimir-Zentralit\"at).} Der Laplace-Operator $\Delop$ ist
das Bild des Casimir-Elements der universellen einh\"ullenden Algebra
$\mathcal{U}(\mathfrak{sl}_2)$. Dieses Element liegt im Zentrum und kommutiert
daher mit der regul\"aren $G$-Wirkung.

\smallskip\noindent
\textbf{Schritt 2 (radiale Faltung).} Der Operator $R_k$ ist ein Mittel von
Rechts-Translationen mit bi-$K$-invariantem Kernel. Da $\Delop$ mit jeder
solchen Translation kommutiert, kommutiert $\Delop$ auch mit ihrem Integral
$R_k$.

\smallskip\noindent
\textbf{Schritt 3 (gemeinsame Diagonalisierung).} \"Aquivalent diagonalisiert die
Laplace-Eigenbasis beide Operatoren:
\[
\Delop\psi_j=(1/4+r_j^2)\psi_j,\qquad R_k\psi_j=\widehat{k}(r_j)\psi_j.
\]
Damit gilt $[\Delop,R_k]=0$ auf dem sph\"arischen Spektralkanal. Folglich ist
$\Delop$ ein nichttrivialer universeller Kommutant f\"ur die relevante
Selberg-Testfunktionsalgebra.
\end{proof}

\begin{remark}[Keine individuelle Phasen-Eigenwert-Aussage]
Die Aussage ist bewusst auf sph\"arische Faltungsoperatoren beschr\"ankt.  Es wird
nicht behauptet, dass ein individueller Geod\"atenfluss-Pullback auf $SM$ eine
allgemeine Laplace-Eigenfunktion $\psi_j$ mit Eigenwert
$e^{im\ell_\gamma r_j}$ multipliziert.  Die korrekte skalare Wirkung ist die
Selberg-/Harish-Chandra-Transformierte $\widehat{k}(r_j)$ eines bi-$K$-invarianten
Kernels im sph\"arischen Kanal.
\end{remark}

\begin{remark}[Warum dies f\"ur Selberg m\"oglich ist, aber nicht f\"ur Riemann]
Das Argument h\"angt von der Existenz eines \emph{sph\"arischen homogenen Kanals}
ab: $G/K$ besitzt eine bi-$K$-invariante Faltungsalgebra, und das Casimir-Element
wirkt dort zentral.  Die Geod\"atensummen der Spurformel werden in genau diesen
Testfunktionskanal integriert. Primzahlen besitzen im getesteten Riemann-Modell
keine vergleichbare kontinuierliche Symmetrie, die eine solche radiale
Faltungsalgebra und einen Casimir-Kommutanten bereitstellt. Dies ist der
pr\"azise Sinn, in dem NE-B den Unterschied sichtbar macht.
\end{remark}

\subsection{Konsequenz f\"ur den numerischen Kommutator-Test}

W\"urde man den numerischen Kommutator-Test aus
\cite[Abschnitt~5.3]{Geiger2026partII} in einem sph\"arischen Selberg-Modell
wiederholen, sollte man einen \emph{gro\ss{}en} Kern f\"ur die
Kommutator-Bedingungsmatrix $C_N$ erhalten, da jedes geeignete Polynom in
$\Delop$ ein Kommutant der radialen Faltungsalgebra ist. Das Spektrum der
Singul\"arwerte w\"urde dann keine Struktur eines isolierten eindimensionalen
Nullraums zeigen, sondern viele kleine Singul\"arwerte.

Dies ist die Diagnostik: \emph{Wenn der v2.0-Kommutator-Test auf einem echten
sph\"arischen Hilbert--P\'olya-Kanal ausgef\"uhrt wird, entdeckt er die
operatorielle Struktur}. Im getesteten Riemann-Trunkierungsmodell entdeckt der
Test nichts Nichttriviales; daraus folgt, dass zumindest diese universelle
Kommutantenroute blockiert ist.

% =============================================================================
\section{Hilbert--P\'olya als Spezialfall von v2.0}
\label{sec:meta}
% =============================================================================

Wir artikulieren nun das Meta-Prinzip, das aus der Kombination der
Riemann- und Selberg-Analysen hervorgeht.

\subsection{Das Hilbert--P\'olya-Programm neu betrachtet}

Die klassische Hilbert--P\'olya-Vermutung schl\"agt vor, die RH zu beweisen, indem ein
selbstadjungierter Operator $H$ auf einem Hilbertraum konstruiert wird, dessen Spektrum die Menge der
Imagin\"arteile $\{\gamma_k\}$ der nicht-trivialen Riemannschen Nullstellen ist. Ein solcher
Operator, sofern er existiert, w\"urde die RH als unmittelbares Korollar liefern
(Selbstadjungiertheit $\Rightarrow$ reelle Eigenwerte $\Rightarrow$ Nullstellen auf
der kritischen Linie).

Einhundert Jahre der Bem\"uhungen haben Kandidaten-Operatoren hervorgebracht:
Berry--Keatings $H = xp$ \cite{BerryKeating1999}, Connes' adelischen
Dirac-Operator \cite{Connes1999} und weitere Ans\"atze.
Doch keiner davon realisiert nachweislich das Riemannsche Spektrum.

\subsection{Die v2.0-Alternative}

Das v2.0-Programm umgeht die Hilbert--P\'olya-Anforderung, indem es
die RH durch eine \emph{Ungleichung quadratischer Formen} (gerade Dominanz
von $\QW_\lambda$) statt durch eine operator-spektrale Identifikation etabliert.
Die Bestandteile sind:
\begin{itemize}
\item Eine Parit\"atszerlegung des umgebenden Hilbertraums.
\item Ein algebraischer \emph{punktweiser} Parit\"atsvorteil (Verschiebungs-Parit\"ats-Lemma).
\item Ein analytischer \emph{Aggregations}-Mechanismus (Frontier-Dominanz).
\item Ein Argument, dass keine einfachere Route gilt (NE-A, NE-B).
\end{itemize}
Entscheidend ist: \emph{es wird kein Operator konstruiert}. Die Route verläuft direkt \"uber die
quadratische Form.

\subsection{Selberg als Kalibrierung: Der sph\"arische operatorielle Kanal existiert}

Theorem~\ref{thm:v2-Sel-RH} zeigt, dass v2.0 die strenge Aussage für die
nicht-trivialen spektralen Selberg-Nullstellen im kritischen Streifen bedingt
auf die C2-/Minimizer-Gates, die gerade Dominanz von
$\QW^\Gamma_\lambda$ und $\lambda_1(\Delop)\ge 1/4$ reduziert.  In diesem Fall existiert im
sph\"arischen Kanal ebenfalls ein Hilbert--P\'olya-/Casimir-Operator $\Delop$
(Proposition~\ref{thm:NEB-fail}), und die klassische Spektraltheorie
best\"atigt unabh\"angig die Schlussfolgerung.  Die beiden Routen koexistieren:
\begin{itemize}
\item Selbergs urspr\"ungliche Spektraltheorie (unbedingte spektrale
  Parametrisierung; strenge kritische Linie unter $\lambda_1(\Delop)\ge 1/4$).
\item Die v2.0-Reduktion (quadratisch: die C2-/Minimizer-Gates und die gerade
  Dominanz von $\QW^\Gamma_\lambda$ implizieren die strenge
  strenge spektrale Nullstellenaussage unter der genannten L\"uckenhypothese).
\end{itemize}
Beide gelangen bei gegebenen Hypothesen auf unterschiedlichen Wegen zur selben
Schlussfolgerung. Die v2.0-Methode ist daher mit der operatorbasierten Route
kompatibel; sie ersetzt die C2-/Minimizer-Identifikation nicht, sondern macht
sie als explizites Gate sichtbar.

\subsection{Riemann als der komplement\"are Fall}

Im Riemann-Setting behauptet NE-B (in der Vermutungs-Form f\"ur den Gesamtraum),
dass der getestete universelle Kommutantenweg keinen nichttrivialen Operator
liefert. Wenn diese Vermutung gilt, ist diese konkrete Hilbert--P\'olya-Route
blockiert. Dies ist der \emph{wissenschaftlich interessante} Fall: v2.0 versucht
eine Beweisroute \"uber die Quadratform, ohne die Existenz genau dieses
Operators vorauszusetzen.

\subsection{Meta-Theorem}

\begin{theorem}[v2.0 Meta-Prinzip, informell]
\label{thm:meta}
Sei $Z$ eine zeta-\"ahnliche Funktion, die eine explizite Formel vom
Riemann--Weil-Typ zulasst (mit reellwertigen logarithmischen Verschiebungen und positiven Gewichten).
Dann gilt:
\begin{enumerate}[label=(\roman*)]
\item Die v2.0-Bestandteile (Verschiebungs-Parit\"ats-Lemma, Frontier-Dominanz,
Parit\"atszerlegung der assoziierten Weilschen quadratischen Form) \"ubertragen sich auf $Z$
mit Modifikationen lediglich in den Dichte- und Gewichtskonstanten.
\item Die v2.0-Route zur entsprechenden Riemannschen Vermutung ist verf\"ugbar,
sobald eine angemessene Frontier-Dichte-Schranke und die n\"otigen
C2-/Minimizer-Gates gelten.
\item Ein Hilbert--P\'olya-Kanal f\"ur $Z$ verlangt, dass das Analogon von NE-B
in der relevanten Operatorenklasse nicht blockiert, d.h.\ die mit $Z$
assoziierten Testfunktionsoperatoren einen nichttrivialen universellen
Kommutanten zulassen.
\item Wenn ein solcher Operator existiert, ist v2.0 mit ihm kompatibel. Wenn der
getestete NE-B-Weg keinen solchen Operator zul\"asst, ist diese spezielle
Kommutantenstrategie blockiert, nicht automatisch jede denkbare
Hilbert--P\'olya-Realisierung.
\end{enumerate}
\end{theorem}

Wir formulieren dies als ein \emph{strukturelles Prinzip} und nicht als ein formales
Theorem, da es davon abh\"angt, was man als zul\"assige Klasse von
zeta-\"ahnlichen Funktionen und als relevante Operatorenklasse betrachtet. F\"ur
klassische L-Funktionen, die eine Standard-Funktionalgleichung und eine
explizite Formel erf\"ullen (Selberg-Klasse, Connes' Familie), wird erwartet,
dass die Quadratform-Komponenten unter den genannten Bedingungen gelten; die
operatorielle Hilbert--P\'olya-Komponente bleibt eine separate Strukturfrage.

\begin{remark}[Verfeinerung im Begleit-Meta-Papier]
\label{rem:meta-paper-refinement}
Das informelle Meta-Prinzip von Theorem~\ref{thm:meta} wurde im
Begleitpapier \cite{GeigerMeta2026} als das \emph{Meta-Theorem} (dortiges Thm.~thm:meta)
durch die Drei-Komponenten-Zerlegung
\[
\mathrm{RH}_X \;=\; \mathrm{GBC}_X \;+\; C2_X \;+\; \mathrm{Hurwitz}_X
\]
technisch pr\"azisiert.
In dieser Terminologie sind die Punkte (iii)--(iv) unseres informellen Theorems~\ref{thm:meta}
die vorsichtige Aussage, dass der Lie-Ursprungs-Zweig einen sph\"arischen
operatoriellen Kanal bereitstellt und der arithmetische Zweig keinen solchen
Kanal durch den getesteten NE-B-Weg liefert.
Der feinere Aspekt, der in \cite{GeigerMeta2026} hervorgehoben wird,
ist die Zerlegung von $C2_X$ in zwei Teilschritte:
\[
C2_X^{\mathrm{parity}}
\quad\text{und}\quad
C2_X^{\mathrm{eigenvec}}.
\]
Der erste bezeichnet den einfachen Grundzustand mit geradem Eigenvektor,
der zweite die Konvergenz der Hermite-Prolate-Approximation. Im aktuellen CORE-Status
ist der Riemannsche $C2/MS2$-Block im CCM/Fourier-Zweig durch die
Drei-Lemma-Mikrocluster-Schlie\ss{}ung des Zookeepers geschlossen. Dies \"andert
nichts an der Beweisverpflichtung f\"ur Selberg: der Laplace--Casimir-Operator
liefert die operatorielle Orientierung, aber die Minimizer-Identifikation bleibt
als explizites Gate im Paper sichtbar. Der Selberg-Fall ist daher eine positive
Kontrolle f\"ur die Normalform und kein abh\"angiges Korollar der Riemannschen
Schlie\ss{}ung.
\end{remark}

\subsection{Umformulierung: Die Weilsche quadratische Form ist fundamental}

Die strukturelle Erkenntnis lautet:

\medskip
\noindent\emph{Die Weilsche quadratische Form $\QW$ ist das fundamentale Objekt; der
Hilbert--P\'olya-Operator ist eine Realisierung, die existieren kann oder auch nicht.}

\medskip
Einhundert Jahre lang wurde die Riemannsche Vermutung durch die
Hilbert--P\'olya-Linse betrachtet: Finde den Operator, beweise die RH. Der
v2.0-Rahmen kehrt diese Sichtweise um: Beweise die passenden Quadratform- und
C2-Gates, und untersuche gesondert, ob ein Hilbert--P\'olya-Operator existiert.
Im Selberg-Fall existiert ein sph\"arischer Casimir-Kanal; im getesteten
Riemann-Modell liefert NE-B keinen vergleichbaren Kommutanten.
Beide F\"alle werden von demselben zugrundeliegenden Rahmenwerk erfasst.

% =============================================================================
\section{Anwendungen und Ausblick}
\label{sec:applications}
% =============================================================================

\subsection{Verallgemeinerte L-Funktionen}

Die v2.0-Methode ist im Prinzip auf jede zeta-\"ahnliche Funktion mit einer
expliziten Formel vom Riemann--Weil-Typ anwendbar. Wir erwarten Folgendes:

\begin{itemize}
\item \emph{Dirichletsche L-Funktionen} $L(s, \chi)$ f\"ur einen Dirichlet-Charakter
$\chi$: die explizite Formel ist die Primsummation gewichtet mit $\chi(p)$.
Das Verschiebungs-Parit\"ats-Lemma l\"asst sich \"ubertragen (der Charakter f\"uhrt eine Phase ein, keinen
Vorzeichenwechsel der geometrischen Parit\"atsstruktur). Frontier-Dominanz \"ubertr\"agt sich
(mit PNT ersetzt durch den PNT f\"ur L-Funktionen, Landau--Siegel-Schranken).
F\"ur den getesteten NE-B-Kommutantenweg erwarten wir weiterhin ein Hindernis,
da $\chi$ keine neue kontinuierliche sph\"arische Symmetrie einf\"uhrt. In der aktuellen Zeta-Zoo-Kartographie ist Dirichlet
nicht einfach dadurch geschlossen, dass man das ungetwistete Riemann-Paket importiert: es bleibt
der Folge-Zweig Weg-D / $\chi$-getwistetes CCM. Eine konkrete Dirichlet-Ausarbeitung
w\"urde die Vorhersagen des Meta-Rahmens operativ
\"uberpr\"ufbar machen --- es ist der nat\"urliche n\"achste Schritt nach der vorliegenden Selberg-Konsistenzpr\"ufung.

\item \emph{Dedekindsche Zetafunktionen} $\zeta_K(s)$ f\"ur einen Zahlk\"orper $K$: die
explizite Formel summiert \"uber Primideale $\mathfrak{p}$ mit dem Gewicht
$(\log N\mathfrak{p}) \cdot N\mathfrak{p}^{-m/2}$. Das Verschiebungs-Parit\"ats-Lemma und der
Frontier-Mechanismus \"ubertragen sich. F\"ur den NE-B-Kommutantenweg erwarten
wir in generischen K\"orpern weiterhin Blockade; zus\"atzliche arithmetische
Symmetrien (etwa CM-Ph\"anomene) m\"ussen separat analysiert werden.

\item \emph{Automorphe L-Funktionen} (Hecke, $\mathrm{GL}(n)$-Spitzenformen):
die explizite Formel l\"asst sich \"ubertragen. Ein Hilbert--P\'olya-Kanal ist
plausibel, wenn die L-Funktion eine konkrete spektrale Interpretation \"uber
einen Laplace-/Casimir-Operator besitzt (zum Beispiel bei geometrisch
realisierten Maass-Formen). Andernfalls bleibt der NE-B-Kommutantenweg eine
separate Pr\"ufung, kein automatisches Ja/Nein-Kriterium.
\end{itemize}

\subsection{Ruelle-Zeta f\"ur Axiom-A-Fl\"usse}

F\"ur einen Axiom-A-Fluss auf einer kompakten Riemannschen Mannigfaltigkeit (ein hyperbolisches dynamisches
System) wird die Ruelle-Zetafunktion definiert durch
\begin{equation}
\zeta_R(s) = \prod_\tau \frac{1}{1 - e^{-s\ell_\tau}},
\end{equation}
wobei das Produkt \"uber primitive periodische Orbits $\tau$ mit minimaler
Periode $\ell_\tau$ gebildet wird. F\"ur glatte Anosov-Fl\"usse liefern Arbeiten
wie \cite{Dolgopyat1998} und verwandte Entwicklungen starke Aussagen zu
analytischer Fortsetzung, Resonanzen und Korrelationszerfall; sie sind jedoch
kein allgemeiner kritische-Linien-Satz im Selberg-Sinn. F\"ur
nicht-hyperbolische Fl\"usse ist eine entsprechende RH-artige Formulierung erst
recht eine Vermutung.

\begin{conjecture}[v2.0 Ruelle-Testproblem]
\label{conj:Ruelle}
F\"ur einen hyperbolischen Fluss mit topologischer Entropie $h$ ist die v2.0-Methode auf
die Ruelle--Weil quadratische Form $\QW^\mathrm{R}_\lambda$ testbar. Zu pr\"ufen
sind insbesondere die Frontier-Dominanz, die Remainder-Kontrolle und die
C2-/Minimizer-Identifikationsgates f\"ur eine m\"ogliche kritische-Linien-Aussage
$\mathrm{Re}(s)=h/2$. In generischen Settings ohne symmetrischen Casimir-Kanal
erwarten wir, dass der NE-B-Kommutantenweg blockiert ist.
\end{conjecture}

Dieses Testproblem ist schwieriger als der Selberg-Kalibrierungsfall, weil es
sich nicht auf eine zugrundeliegende Riemannsche symmetrische Struktur verlassen
kann. F\"ur generische Axiom-A-Fl\"usse ist kein Laplace-/Casimir-Operator
verf\"ugbar; die v2.0-Methode w\"are daher, wenn die Gates gelingen, eine
Quadratformroute statt eines bereits vorhandenen Spektralbeweises.

\subsection{Ihara-Zeta und die diskret-kontinuierliche Achse}
\label{sec:ihara}

Die Ihara-Zetafunktion $\zeta_\Gamma(s)$ eines endlichen zusammenh\"angenden $(q+1)$-regul\"aren
Graphen $\Gamma$ ist definiert durch
\begin{equation}
\zeta_\Gamma(s)^{-1}
  = (1-s^2)^{r-1}\det\!\bigl(I - sA + qs^2 I\bigr),
\end{equation}
wobei $A$ die Adjazenzmatrix und $r$ der Zyklusrang von $\Gamma$ ist
\cite{Ihara1966,Hashimoto1989,Bass1992}.  Die Determinantenformel
wird manchmal als Ihara--Bass-Formel bezeichnet; die Beweise von Hashimoto (1989)
und Bass (1992) etablierten sie in voller Allgemeinheit.  Die Ihara-RH ist die Aussage, dass alle
nicht-trivialen Pole auf dem Kreis $|s| = q^{-1/2}$ liegen, was \"aquivalent dazu ist, dass
$\Gamma$ ein Ramanujan-Graph ist.  Dies wurde f\"ur explizite
Ramanujan-Familien bewiesen, die aus der $\mathrm{PGL}(2)$-Arithmetik konstruiert sind
\cite{LPS1988}.

\medskip
Im v2.0 / Zookeeper-Rahmen ist der Ihara-Fall \emph{diskret-operatorisch}
(SGE-YES), parallel zum \emph{kontinuierlich-operatorischen} Selberg-Fall.
Der Graph-Laplace-Operator $L = (q+1)I - A$ kommutiert mit allen Automorphismen von $\Gamma$
und damit mit der Kanten-Adjazenz-Transfermatrix $T$.  Dies macht $L$ zum
diskreten Analogon des hyperbolischen Laplace-Operators $\Delta_\Gamma$: ein universeller
Kommutant, der NE-B f\"ur Ihara aus exakt demselben strukturellen Grund wie im Selberg-Fall inhaltlos macht.

Die drei-Wege-Vergleichsachse im Zookeeper-Rahmen lautet:

\begin{center}
\scriptsize
\setlength{\tabcolsep}{2pt}
\begin{tabular}{@{}L{2.0cm}L{1.7cm}L{5.2cm}L{2.0cm}@{}}
\toprule
\shortstack[l]{\textbf{Zeta-}\\\textbf{funktion}} & \textbf{Typ} & \textbf{Kommutant} & \shortstack[l]{\textbf{v2.0-}\\\textbf{Klasse}} \\
\midrule
\shortstack[l]{Selberg\\$Z_\Gamma(s)$} & \shortstack[l]{kont.\\geom.} & \shortstack[l]{$\Delta_\Gamma$\\(Laplace--Beltrami)} & SGE-YES \\
Ihara $\zeta_\Gamma(s)$ & \shortstack[l]{diskret,\\ komb.} & $L = (q{+}1)I - A$ (Graph-Laplace) & SGE-YES \\
Riemann $\zeta(s)$ & arithmetisch & keiner (NE-B gilt) & SGE-NO \\
\bottomrule
\end{tabular}
\end{center}

\medskip\noindent
In beiden SGE-YES-F\"allen wird die strenge RH-artige Aussage durch eine klassische
Spektralschranke gesteuert (Selbergs Bedingung ohne Ausnahmeeigenwerte
$\lambda_1(\Delop)\ge 1/4$; Ramanujans $|$Eigenwert$| \le 2\sqrt{q}$),
und der v2.0-Rahmen liefert eine \emph{konsistente bedingte Neuableitung} durch gerade Dominanz
der zugeh\"origen Weilschen quadratischen Form.  F\"ur die Riemannsche Zetafunktion existiert
ein solcher Kommutant nicht; die v2.0-Methode ist die einzige verf\"ugbare strukturelle Route zur RH.

\subsection{Physikalische Interpretation}

Das Selberg-Setting entspricht physikalisch einem Quantenteilchen auf einer
hyperbolischen Fl\"ache (``Quantenchaos''); die Laplace-Eigenwerte sind die
Energieniveaus. Das v2.0 Verschiebungs-Parit\"ats-Lemma besagt, dass die Parit\"at der
Eigenfunktionen an der Schwelle hin zu geraden Funktionen verzerrt ist (biased). Dies ist eine
Aussage \"uber die Parit\"atsklassen-Verteilung von Quanteneigenzust\"anden.

F\"ur Axiom-A-Fl\"usse (Billards, chaotische Hohlr\"aume) prognostiziert der v2.0-Rahmen
GUE-artige Statistiken f\"ur die Ruelle-Nullstellen, wie in der BGS (Bohigas--Giannoni--Schmit)
Vermutung \cite{BGS1984}; siehe auch \cite{RudnickSarnak1996} f\"ur verwandte GUE-Statistiken
von Nullstellen von Haupt-$L$-Funktionen. Die v2.0-Methodik k\"onnte eine
strukturelle Erkl\"arung dieser Vermutung bieten.

% =============================================================================
\section{Zookeeper-Konsistenz: Selberg als SGE-YES Test}
\label{sec:zookeeper-selberg}

Der Selberg-Fall ist auch ein n\"utzlicher Kalibrierungspunkt f\"ur die Post-Zookeeper-Version
des Zeta-Zoos.  In dieser Sprache ist Selberg ein SGE-YES-System:
der relevante Operator existiert klassisch, und der v2.0-Transfer muss
daher einen bekannten positiven Fall reproduzieren anstatt eine neue
Aussage von Grund auf zu kreieren.

Tabelle~\ref{tab:transfer-slots} erfasst die f\"unf Transfer-Slots und
ihre Belegungen in den Selberg- und Riemann-Settings.

\begin{table}[htp]
\centering
\footnotesize
\setlength{\tabcolsep}{4pt}
\caption{Die f\"unf SGE-YES Transfer-Slots f\"ur Selberg vs.\ Riemann.
  Ein H\"akchen (\checkmark) zeigt an, dass der Slot klassisch besetzt ist;
  ``v2.0'' zeigt an, dass er durch das v2.0-Rahmenwerk etabliert wird; ``hyp.''
  markiert die zus\"atzliche Spektrall\"ucken-Hypothese f\"ur die strenge Linienform.}
\label{tab:transfer-slots}
\begin{tabularx}{\textwidth}{@{}L{2.3cm}L{4.95cm}L{4.85cm}@{}}
\toprule
\textbf{Slot} & \textbf{Selberg (SGE-YES)} & \textbf{Riemann} \\
\midrule
Operator / Form
  & $\Delop$ (klassischer Casimir; kommutiert mit sph\"arischen radialen Faltungen; Intervall-Shift-Intertwiner nicht konstruiert) \checkmark
  & $\QW_\lambda$ (kein klassischer HP-Operator; getesteter NE-B-Weg blockiert) \\[4pt]
Defekt
  & Casimir-Residual $S_\Gamma(L,I)=r_\Gamma/g_I$ (Companion-Refinement)
  & L\"ucke zur geraden Dominanz (vermutet $\Rightarrow$ etabliert durch v2.0) \\[4pt]
Approximation
  & L\"angen-Cutoffs, trunkierte Spurformel \checkmark
  & Galerkin-Trunkierung $D_N(r)$ (v2.0) \\[4pt]
L\"ucke
  & Keine-Ausnahmeeigenwert-L\"ucke $\lambda_1(\Delop)\ge \tfrac{1}{4}$ (hyp.)
  & Gerade-Dominanz-L\"ucke $\Delta(\lambda)>0$ (v2.0) \\[4pt]
Grenzwert
  & Volle Selberg-Spurformel \cite{Selberg1956} \checkmark
  & Volle Weilsche explizite Formel \\
\bottomrule
\end{tabularx}
\end{table}

Die Defekt-Zeile verwendet die gesch\"arfte Beweisnotiz-Formulierung:
Der Selberg-Defekt soll als Residual-zu-L\"ucke-Signal
$S_\Gamma(L,I)=r_\Gamma(L,I)/g_I$ gemessen werden, wobei $g_I$ die \"au\ss{}ere
Casimir-Spektrall\"ucke ist.  Dies ist ein Companion-Refinement des positiven
Kontrollfalls, kein zus\"atzliches unbedingtes Selberg-RH-Theorem.

\begin{remark}[Galerkin vs.\ spektrale Trunkierung]
\label{rem:galerkin-vs-spectral}
Im Riemann-Fall ist die Approximation ein \emph{Galerkin-Cutoff}: die
unendlichdimensionale quadratische Form $\QW_\lambda$ wird durch die
endliche Matrix $D_N(r)$ \"uber eine Kosinus/Sinus-Basis der Gr\"o\ss{}e $N$ approximiert.  F\"ur Selberg
ist die Approximation stattdessen ein \emph{spektraler Cutoff}: die geod\"atische Summe wird
bei der L\"ange $\ell_\gamma \le 2L$ trunkiert, und Bemerkung~\ref{rem:truncation-exact}
zeigt, dass diese Trunkierung exakt und nicht approximativ ist.  Der Zookeeper-\emph{Approximations}-Slot
verwendet daher in beiden F\"allen unterschiedliche, aber analoge Technologien.
\end{remark}

Somit spielt Selberg eine diagnostische Rolle.  W\"urde die Zookeeper-Normalform hier keine
positive L\"ucke detektieren, w\"are das Mapping falsch.  Umgekehrt
erkl\"art die Tatsache, dass der Laplace-/Casimir-Operator mit den
sph\"arischen radialen Faltungsoperatoren kommutiert, warum das relevante
NE-B-Hindernis im Selberg-Setting fehlt und hebt den Kontrast zum Riemann-Fall
hervor: f\"ur Selberg ist ein sph\"arischer Hilbert--P\'olya-Kanal pr\"asent; f\"ur
Riemann ist die urspr\"ungliche universelle Kommutantenroute strukturell
blockiert, w\"ahrend das separate Zookeeper-Paket den Riemannschen $C2/MS2$-Kern
im CCM/Fourier-Zweig schlie\ss{}t.  Selberg kalibriert folglich die
SGE-YES-Seite des Randtheorems als positive Kontrolle; es ist kein
zus\"atzlicher offener Blocker und kein Beweis daf\"ur, dass das
Prime-Hub/OP5-Graphenproblem gel\"ost wurde.

\begin{proposition}[Selberg--Zookeeper Konsistenz]
\label{prop:selberg-zookeeper}
Das v2.0-Rahmenwerk ordnet die Selbergsche Zetafunktion
$\Sel(s)$ auf der genannten Hypothesenebene als positiven
SGE-YES-Kontrollfall ein.  Spezifisch:
\begin{enumerate}[label=(\roman*)]
  \item Die C2-/Minimizer-Gates f\"ur $\QW^\Gamma_\lambda$ implizieren zusammen
    mit $\lambda_1(\Delop)\ge 1/4$ die strenge Aussage für die nicht-trivialen
    spektralen Nullstellen im kritischen Streifen (Theorem~\ref{thm:v2-Sel-RH}).
  \item Der Laplace--Beltrami-/Casimir-Operator $\Delop$ ist ein universeller
    Kommutant f\"ur den sph\"arischen radialen Faltungskanal, weshalb das
    relevante NE-B-Hindernis f\"ur Selberg fehlt
    (Proposition~\ref{thm:NEB-fail}).
  \item Die Zookeeper-Transfer-Slots sind identifiziert; der L\"ucken-Slot ist
    bedingt auf die Hypothese ohne Ausnahmeeigenwerte
    (Tabelle~\ref{tab:transfer-slots}).
  \item Die v2.0-Klassifizierung ist mit Selbergs klassischer Spektraltheorie
    \"uber $\Delop$ auf derselben Hypothesenebene konsistent.
\end{enumerate}
\end{proposition}

\section{Fazit}
\label{sec:conclusion}
% =============================================================================

Wir haben die v2.0-Methodik an der Selbergschen Zetafunktion getestet. Das
Verschiebungs-Parit\"ats-Lemma, die Frontier-Dominanz, die Weilsche quadratische Form
und NE-A \"ubertragen sich auf das Selberg-Setting mit bescheidenen Modifikationen
(haupts\"achlich Konstanten und Dichteregime), wobei die Frontier-Aussage als
bedingte Dominanzschranke zu lesen ist. Die C2-/Minimizer-Gates f\"ur die
Selberg--Weilsche quadratische Form $\QW^\Gamma_\lambda$ liefern zusammen mit
$\lambda_1(\Delop)\ge 1/4$ eine bedingte v2.0-Reduktion der strengen Aussage
für die nicht-trivialen spektralen Nullstellen im kritischen Streifen,
konsistent mit der klassischen spektralen Beschreibung.

Die entscheidende strukturelle Erkenntnis ist, dass das relevante
\emph{NE-B-Hindernis im sph\"arischen Selberg-Kanal fehlt}: Der
Laplace--Beltrami-/Casimir-Operator $\Delop$ ist ein universeller Kommutant
f\"ur die radiale sph\"arische Faltungsalgebra. Dies ist das Kennzeichen eines
operatoriellen Hilbert--P\'olya-Kanals --- und der v2.0-Kommutator-Test,
ausgef\"uhrt auf einem solchen Kanal, w\"urde ihn erkennen. Es ist dagegen keine
Behauptung, dass einzelne Geod\"atenfluss-Pullbacks auf $SM$ bereits mit den
trunkierten Intervall-Shifts identifiziert wurden.

Wenn wir die Riemann- und Selberg-Bilder kombinieren, formulieren wir ein Meta-Prinzip:
Das v2.0-Rahmenwerk ist ein gemeinsamer Quadratform-Test, und das
Hilbert--P\'olya-Programm ist ein operatorieller Spezialfall, der zus\"atzliche
Kommutantenstruktur verlangt. Im Selberg-Fall liegt diese Struktur im
sph\"arischen Casimir-Kanal vor. Im Riemann-Fall blockiert NE-B zumindest den
getesteten universellen Kommutantenweg; v2.0 muss daher durch die Quadratform-
und C2-Gates laufen.

\medskip
\noindent\textbf{Prinzip.} \emph{Die Weilsche quadratische Form ist das
gemeinsame Testobjekt. Der Operator ist zus\"atzliche Struktur.}

\medskip
Ein Jahrhundert des Nachdenkens hat den Operator als das Herz der
Riemannschen Vermutung identifiziert. Der v2.0-Rahmen legt stattdessen nahe,
dass die Ungleichung quadratischer Formen das gemeinsame Herzst\"uck ist und der
Operator, wo er existiert, eine zus\"atzliche geometrisch-symmetrische
Realisierung darstellt.


% =============================================================================
\section*{KI-Offenlegung}
\addcontentsline{toc}{section}{KI-Offenlegung}
% =============================================================================

\noindent\textbf{Erklärung zum Einsatz von KI-Werkzeugen}

\smallskip
\small

\begin{sloppypar}
\noindent
Dieses Manuskript wurde unter umfassender Unterstützung durch KI-gestützte
Sprachmodelle (Claude, Anthropic; Gemini, Google; ChatGPT, OpenAI) erstellt.
Die KI-Werkzeuge wurden in mehreren Phasen des Forschungs- und
Schreibprozesses eingesetzt, darunter konzeptionelle Diskussion,
Literaturrecherche, Strukturierung, Formulierungshilfe, Analyse- und
Synthesevorschläge, kritische Gegenprüfung, sprachliche Überarbeitung,
Übersetzung, Zitations- und Quellenarbeit sowie technische
\LaTeX-Formatierung, Satzvorbereitung und Kompilierungsunterstützung.
\end{sloppypar}

\medskip

\noindent
Für die rechnerischen Bestandteile dieser Arbeit --- insbesondere die
C2c-Fensterledger-Skripte und ihre Auswertungsläufe
(\Cref{sec:c2c-real-ledger}) --- wurden KI-Werkzeuge zusätzlich zur
Erstellung, Anpassung, Umschreibung, Dokumentation und Prüfung von
Python-Skripten sowie zur Planung, Durchführung und Betreuung von
Berechnungsläufen eingesetzt. Dies umfasste insbesondere die
Weiterentwicklung bestehender Skripte, die Vorbereitung nachvollziehbarer
Ausgaben, die Auswertung von Zwischenergebnissen, Reanalysen und die
Dokumentation von Tests, Verfahren und Ergebnissen.

\medskip

\noindent
Zur Qualitätssicherung wurden mehrere komplementäre Verfahren eingesetzt:
Modell-Cross-Checks, mehrstufige Review-Pipelines, explizit
widerlegungsorientierte Prüfungen, Quellen- und Zitationschecks,
sprachliche Gegenlesungen, Reanalysen, reproduzierbare Skripte, gespeicherte
Resultatdateien, Testläufe, Hash-Vergleiche, erneute Ausführungen sowie eine
ordnerbasierte Forschungspipeline mit Projekt-Guidelines, vorgegebenen
Workflows, Dokumentationsregeln und dokumentierten Korrekturschritten.

\medskip

\noindent\textit{Modell- und Workflow-Zuordnung.}

\medskip

\begin{sloppypar}
\noindent
In der bisherigen Forschungsarbeit wurden einfachere Aufgaben typischerweise
mit Claude Opus/Sonnet~4.x, GPT-5.4 und Gemini~3.x~Pro bearbeitet; Gemini
wurde insbesondere auch für Übersetzungen eingesetzt. Anspruchsvollere
Aufgaben wurden mit Claude Opus~4.6/4.7, Claude Fable~5 und GPT-5.5
durchgeführt. Claude Opus wurde häufig als dialogischer Forschungspartner
zur Entwicklung und Fortführung des Forschungsprozesses eingesetzt, während
GPT-5.5 (über die Codex-Werkzeugkette) verstärkt für Reviews, Audits,
Gegenprüfungen und kritische Nachanalysen genutzt wurde --- einschließlich
des adversariellen Gegenreviews des Real-Daten-Fensterledgers aus
\Cref{sec:c2c-real-ledger}. Der Real-Daten-Ledger-Lauf und seine
Integration in die vorliegende Fassung erfolgten mit Claude Fable~5. Alle
genannten Modellfamilien wurden, soweit einschlägig, auch für \LaTeX-Satz,
technische Formatierung, Kompilierung und Fehlerdiagnose eingesetzt.
Methodisch kamen iterative Schleifen, Modell-Cross-Checks,
Interchat-Gespräche zwischen Claude-, Gemini- und GPT-basierten Systemen
sowie ein mehrstufiges 7-Schritte-Review-Verfahren mit explizitem
Widerlegungsagenten und anschließendem Heilungs- bzw.\ Reparaturversuch zum
Einsatz. Der Mensch fungierte als Gedächtnis, verbindendes und antreibendes
Element, reflektierende Instanz und Orchestrator des Gesamtprozesses; er
brachte zudem häufig Lösungen gerade durch den Versuch ein, die offenen
Probleme selbst zu verstehen.
\end{sloppypar}

\medskip

\noindent
Der Autor bestimmte Forschungsfrage, methodische Ausrichtung, zentrale
Thesen, Bewertung der Quellen, fachliche Interpretation und
Schlussfolgerungen. Die volle wissenschaftliche und inhaltliche
Verantwortung für diesen Artikel liegt beim Autor, Lukas Geiger.

\normalsize

% =============================================================================
% Bibliographie
% =============================================================================

\begingroup
\small
\begingroup
\small
\begin{thebibliography}{99}
\setlength{\itemsep}{0pt}
\setlength{\parskip}{0pt}
\setlength{\itemsep}{0.15em}
\setlength{\parsep}{0pt}
\setlength{\parskip}{0pt}
\bibitem{Bass1992}
H.\ Bass,
\emph{The Ihara--Selberg zeta function of a tree lattice},
Int.\ J.\ Math.\ \textbf{3} (1992), no.~6, 717--797.
\href{https://doi.org/10.1142/S0129167X92000357}{doi:10.1142/S0129167X92000357}.

\bibitem{BerryKeating1999}
M. V. Berry und J. P. Keating,
\emph{The Riemann zeros and eigenvalue asymptotics},
SIAM Review \textbf{41} (1999), 236--266.
\href{https://doi.org/10.1137/S0036144598347497}{doi:10.1137/S0036144598347497}.

\bibitem{BGS1984}
O. Bohigas, M.-J. Giannoni und C. Schmit,
\emph{Characterization of chaotic quantum spectra and universality of level fluctuation laws},
Phys. Rev. Lett. \textbf{52} (1984), 1--4.
\href{https://doi.org/10.1103/PhysRevLett.52.1}{doi:10.1103/PhysRevLett.52.1}.

\bibitem{BostConnes1995}
J.-B. Bost und A. Connes,
\emph{Hecke algebras, type III factors and phase transitions with spontaneous symmetry breaking in number theory},
Selecta Math. (N.S.) \textbf{1} (1995), 411--457.
\href{https://doi.org/10.1007/BF01589495}{doi:10.1007/BF01589495}.

\bibitem{Connes1999}
A. Connes,
\emph{Trace formula in noncommutative geometry and the zeros of the Riemann zeta function},
Selecta Math. (N.S.) \textbf{5} (1999), 29--106.
\href{https://doi.org/10.1007/s000290050042}{doi:10.1007/s000290050042}.

\bibitem{Connes2026}
A. Connes,
\emph{The Riemann Hypothesis: Past, Present and a Letter Through Time},
arXiv:2602.04022 (Februar 2026).
\href{https://arxiv.org/abs/2602.04022}{arXiv:2602.04022}.

\bibitem{ConnesMoscovici2022}
A. Connes und H. Moscovici,
\emph{The UV prolate spectrum matches the zeros of zeta},
Proc. Natl. Acad. Sci. USA \textbf{119} (2022), no.~22, e2123174119.
\href{https://doi.org/10.1073/pnas.2123174119}{doi:10.1073/pnas.2123174119}.
Auch: arXiv:2112.05500.

\bibitem{Deninger1994}
C. Deninger,
\emph{Motivic $L$-functions and regularized determinants},
in \emph{Motives} (Seattle, WA, 1991),
Proc. Sympos. Pure Math. \textbf{55}, Part~1, Amer. Math. Soc., Providence, RI, 1994, pp.~707--743.
\href{https://doi.org/10.1090/pspum/055.1/1265547}{doi:10.1090/pspum/055.1/1265547}.

\bibitem{Dolgopyat1998}
D. Dolgopyat,
\emph{On decay of correlations in Anosov flows},
Ann. of Math. (2) \textbf{147} (1998), 357--390.
\href{https://doi.org/10.2307/121012}{doi:10.2307/121012}.

\bibitem{Geiger2026partII}
L. Geiger,
\emph{From Landscape to Proof: The Even-Dominance Route to the Riemann Hypothesis},
Zenodo-Preprint (2026), v2.4. \href{https://doi.org/10.5281/zenodo.20358728}{doi:10.5281/zenodo.20358728}.

\bibitem{GeigerMeta2026}
L. Geiger,
\emph{The Meta-Galerkin Bridge: Explicit Formula Principle and the Three-Component Decomposition of Riemann-Hypothesis-Type Statements},
Begleitpapier (2026, in Vorbereitung).

\bibitem{Hashimoto1989}
K.\ Hashimoto,
\emph{Zeta functions of finite graphs and representations of $p$-adic groups},
in \emph{Automorphic Forms and Geometry of Arithmetic Varieties},
Adv.\ Stud.\ Pure Math.\ \textbf{15}, Academic Press, 1989, pp.~211--280.
\href{https://doi.org/10.1016/B978-0-12-330580-0.50015-X}{doi:10.1016/B978-0-12-330580-0.50015-X}.

\bibitem{Hejhal1976}
D. A. Hejhal,
\emph{The Selberg Trace Formula for $\mathrm{PSL}(2,\RR)$}, Vol.~1,
Lecture Notes in Mathematics \textbf{548}, Springer-Verlag, 1976.

\bibitem{Huber1959}
H. Huber,
\emph{Zur analytischen Theorie hyperbolischer Raumformen und Bewegungsgruppen},
Math. Ann. \textbf{138} (1959), 1--26.
\href{https://doi.org/10.1007/BF01369663}{doi:10.1007/BF01369663}.

\bibitem{Ihara1966}
Y. Ihara,
\emph{On discrete subgroups of the two by two projective linear group over $p$-adic fields},
J. Math. Soc. Japan \textbf{18} (1966), no.~3, 219--235.
\href{https://doi.org/10.2969/jmsj/01830219}{doi:10.2969/jmsj/01830219}.

\bibitem{Iwaniec2002}
H. Iwaniec,
\emph{Spectral Methods of Automorphic Forms},
2nd ed., Graduate Studies in Mathematics \textbf{53}, Amer. Math. Soc., 2002.

\bibitem{LowryDuda2025}
D. Lowry-Duda,
\emph{A database of rigorous Maass forms},
arXiv:2502.01442 (2025).
\href{https://arxiv.org/abs/2502.01442}{arXiv:2502.01442}.

\bibitem{LPS1988}
A. Lubotzky, R. Phillips und P. Sarnak,
\emph{Ramanujan graphs},
Combinatorica \textbf{8} (1988), no.~3, 261--277.
\href{https://doi.org/10.1007/BF02126799}{doi:10.1007/BF02126799}.

\bibitem{RudnickSarnak1996}
Z. Rudnick und P. Sarnak,
\emph{Zeros of principal $L$-functions and random matrix theory},
Duke Math. J. \textbf{81} (1996), no.~2, 269--322.
\href{https://doi.org/10.1215/S0012-7094-96-08115-6}{doi:10.1215/S0012-7094-96-08115-6}.

\bibitem{Sarnak1981}
P. Sarnak,
\emph{Asymptotic behavior of periodic orbits of the horocycle flow and Eisenstein series},
Comm. Pure Appl. Math. \textbf{34} (1981), 719--739.
\href{https://doi.org/10.1002/cpa.3160340602}{doi:10.1002/cpa.3160340602}.

\bibitem{Selberg1956}
A. Selberg,
\emph{Harmonic analysis and discontinuous groups in weakly symmetric Riemannian spaces with applications to Dirichlet series},
J. Indian Math. Soc. \textbf{20} (1956), 47--87.

\bibitem{Selberg1965}
A. Selberg,
\emph{On the estimation of Fourier coefficients of modular forms},
Proc. Sympos. Pure Math. \textbf{VIII}, Amer. Math. Soc., Providence, RI, 1965, pp.~1--15.

\bibitem{SeymourHowell2022}
A. Seymour-Howell,
\emph{Rigorous computation of Maass cusp forms of squarefree level},
Res. Number Theory \textbf{8} (2022), no.~4, Paper No.~83.
\href{https://doi.org/10.1007/s40993-022-00393-y}{doi:10.1007/s40993-022-00393-y}.

\bibitem{StrohmaierUski2013}
A. Strohmaier und V. Uski,
\emph{An algorithm for the computation of eigenvalues, spectral zeta functions and zeta-determinants on hyperbolic surfaces},
Comm. Math. Phys. \textbf{317} (2013), 827--869.
\href{https://doi.org/10.1007/s00220-012-1557-1}{doi:10.1007/s00220-012-1557-1}.

\bibitem{Weil1952}
A. Weil,
\emph{Sur les ``formules explicites'' de la th\'eorie des nombres premiers},
Comm.\ S\'em.\ Math.\ Univ.\ Lund (d\'edi\'e \`a M.\ Riesz) (1952), 252--265.
\end{thebibliography}
\endgroup
\endgroup

\end{document}
